\chapter{Viscous stress and implicit solve}\label{ch:viscosity}
\chapterstatus[bc0da10fc3b7ff326073edc407ec29ff4a22562e]{reference}
Viscosity is the one implicit substep of the split time step: after the
transport and the body acceleration, before the phase forces and the
projection (\cref{eq:forces:stages}), the velocity is advanced by the
implicit midpoint rule. Two operators implement it. A single phase has the
constant kinematic viscosity $\visc/\dens$ and a componentwise staggered
Laplacian on the dual graph of \cref{ch:momentum-transport}, with geometric
dual-face contacts and a symmetric dissipative stiffness at level transitions;
a cut-cell variant adds the embedded walls. Two phases have the
coupled variable stress $\divergence(2\visc\strainrate)$, assembled as
$\strainmat^{\transpose}\stressw\strainmat$ from normal-stress rows in the
cells, shear rows at the cell corners and wall rows, with a mixture of the
phase viscosities that is harmonic (series) by default and geometric as an
option; the two-velocity model diffuses the mass-weighted mixture velocity with
the same operator and gives both phases the same increment (with the mimetic
projection: the same stress in its metric), so the rule applies
to both velocity models, and then relaxes the phase-velocity jump of the
mixed faces by an implicit Navier-slip interphase drag. This chapter
states the continuous model, the semi-discrete operators with their walls,
the fully discrete midpoint solves with their stopping criteria, and the
exact identities that hold for them: dissipation of the mass-weighted
kinetic energy, conservation of momentum on periodic meshes, the
unchanged velocity jump of the two-velocity mixture step and the momentum
conservation and dissipation of the drag. Derivations are in the
thesis~\thesis{\S7.5}; the linear solvers and their multigrid
preconditioners are in \cref{ch:solvers}.

\section{Continuous model}\label{sec:viscosity:continuous}
The viscous part of the phase momentum
balance~\cref{eq:forces:cont-momentum} is
\begin{equation}\label{eq:viscosity:cont}
  \dens^{\phase}\partial_\timevar\Velo^{\phase}=\divergence\roundpar{2\visc^{\phase}\strainrate(\Velo^{\phase})}+\dots,
  \qquad
  \strainrate(\Velo)=\half\roundpar{\gradient\Velo+(\gradient\Velo)^{\transpose}},
\end{equation}
which for a constant viscosity and $\divergence\Velo=0$ is $\visc\Delta\Velo$.
The one-velocity model has a single velocity with the piecewise constant
$\dens$ and $\visc$~\thesis{eq.~(2.9)}; the thesis replaces the sharp
viscosity by the volume-fraction-weighted geometric mean of the phase
viscosities, a diffuse model chosen so that the one- and two-velocity
formulations share it~\thesis{\S7.5.1}. The code implements that rule and, as
the default, the volume-fraction-weighted harmonic mean, the series
conductance across the interface that a centred shear difference needs;
\cref{sec:viscosity:stress} states both and the deviation. With stationary no-slip or free-slip
walls, or periodic boundaries, the viscous term dissipates kinetic energy,
\begin{equation}\label{eq:viscosity:dissipation}
  \frac{\dd}{\dd\timevar}\integral{\domain^{\open}}{\tfrac12\dens\abs{\Velo}^2}{V}
  =-\integral{\domain^{\open}}{2\visc\,\strainrate(\Velo):\strainrate(\Velo)}{V}\le0
  \qquad(\text{viscous part only}),
\end{equation}
and conserves momentum $\integral{\domain}{\dens\Velo}{V}$ on a periodic
domain. The wall conditions are the stationary or moving no-slip
condition $\Velo=\Velo^{\physbnd}$, the free-slip condition
$\Velo\cdot\normal=0$ with zero tangential stress, a prescribed inflow
velocity and, at outflow sides, a zero normal derivative.

\section{Single-phase operator}\label{sec:viscosity:single}
The unknowns of component $k$ are the face velocities $\velo_f$,
$f\in\facesaxis{k}\cap\facesint$, with the transport volumes $\stagvol{f}$
of~\cref{eq:momentum:volumes} as the diagonal of $\dualvolmat$. The
operator is built from the dual faces of the momentum schedule of
\cref{sec:momentum:dual}. For connections $g$ along the component axis $t=k$,
the net measure and normal distance are
\begin{equation}\label{eq:viscosity:connections}
  \meas{g}=\Big|\sum_{g'}\sum_{f'\in\facesaxis{t}}\dweight{g'}{f'}\meas{f'}\Big|,
  \qquad
  \conndist{g}=\abs{(\pos{\dplus{g}}-\pos{\dminus{g}})_t},
\end{equation}
summed over the dual faces $g'$ of the pair; the nearest periodic image is
used and negligible net measures are dropped. Transverse contacts $t\ne k$
are rebuilt from overlaps of bisected cell traces, including the contact
between segments of a split side, as described below for open traces. Without
a solid these are the full traces. This avoids using advection bookkeeping
weights as geometric diffusion measures (\issue{78}, \issue{80}).

Across a connection the directional derivative is a row
\begin{equation}\label{eq:viscosity:gradient}
  \conngrad{g}\avarstag=\sum_f\conncoef{g}{f}\avarstag_f .
\end{equation}
Aligned contacts use a two-point difference. Fine--coarse normal contacts use
an interpolation of the fine siblings when available; transverse contacts
with cancelling weighted offsets share a centred derivative. An offset contact
that is left without such a group is closed against the wall where the wall
value is known and otherwise takes a linear-exact two-point difference on the
tangential level of one of its ends, as described below. With the
wall terms included, the integrated stiffness and Laplacian are
\begin{equation}\label{eq:viscosity:stiffness}
  \viscstiff{k}=\sym{\sum_{g\in\viscconn{k}}\meas{g}\conndist{g}\,
    \conncoef{g}{}^{\transpose}\conncoef{g}{}
    +\sum_{w\in\wallconns{k}}\wallcond{w}\,
      \unitvec{f(w)}\unitvec{f(w)}^{\transpose}},
\end{equation}
\begin{equation}\label{eq:viscosity:laplacian}
  \lapstag{k}=-\dualvolmat^{-1}\viscstiff{k}.
\end{equation}
The symmetrisation removes only round-off. The interior terms are non-negative
weighted squares; the signed wall-contact corrections complete non-negative
wall-aware squares. Hence $\viscstiff{k}$ is symmetric positive semidefinite.
On a periodic mesh constants are in its kernel and the rate conserves
momentum. On a uniform mesh this is the five-point Laplacian.

The previous quadratically exact replacement of individual transition rows
did not preserve the transpose rows, making $\dualvolmat\lapstag{k}$
nonsymmetric and its symmetric part indefinite. The old midpoint step could
increase energy by factors up to $1.1$ (walled) and $2.7$ (periodic) on the
audit fixtures (\code{proof_review/eq-viscosity-laplacian/}). The shared
geometric stiffness removes this growth for every step size (\issue{80}).
Transition rows may have $O(1)$ pointwise truncation error, but the composite
manufactured Helmholtz solution approaches second order, and the observed
slopes depend on where the transition crosses the solution: they are no longer
uniformly above $1.8$. On the audit's walled half-refined mesh, with $16$,
$32$ and $64$ base cells per direction and a sine--sine manufactured
Helmholtz solution, the previous maximum-norm orders were $2.36$, $2.62$ and
the dual-volume $L^2$ orders were $2.61$, $2.67$. The geometric stiffness
gives maximum-norm orders of $1.84$, $1.92$ on a half-refined strip and
$1.78$, $1.88$ on a refined corner, and dual-volume $L^2$ orders of $1.99$,
$2.00$ and $2.01$, $2.01$, respectively (base cells $16\to32\to64$;
\code{test/operators/staggered_diffusion.jl}). Thus the solution remains
approximately second order but the old fixture's superconvergent orders are
not retained.
For stationary walls or periodic boundaries, put
$\mathsf{B}=\dualvolmat^{-1/2}\viscstiff{k}\dualvolmat^{-1/2}$.
Its eigenvalues are non-negative, so for every $a=\dt\visc/\dens\ge0$,
\[
  \norm{\dualvolmat^{1/2}(\idmat-a\lapstag{k})^{-1}\dualvolmat^{-1/2}}_2
  =\norm{(\idmat+a\mathsf{B})^{-1}}_2\le1 .
\]
The midpoint amplification is likewise bounded by one when its explicit
and implicit halves use the same velocity. The split
step~\cref{eq:viscosity:single-midpoint} solves for
$\velo^{\stage{2}}$ from the transported velocity $\velo^{\stage{1}}$ and
evaluates its explicit half at the old velocity $\velo^{\timelevel{n}}$: its
right-hand side is $(\idmat+a\lapstag{k})\velo^{\timelevel{n}}
+(\velo^{\stage{1}}-\velo^{\timelevel{n}})$. With
$\mathsf{M}=\dualvolmat^{1/2}(\idmat-a\lapstag{k})^{-1}(\idmat+a\lapstag{k})\dualvolmat^{-1/2}$,
$\norm{\mathsf{M}}_2\le1$, and, before the wall source,
\begin{equation}\label{eq:viscosity:split-bound}
  \big((\velo^{\stage{2}})^{\transpose}\dualvolmat\,\velo^{\stage{2}}\big)^{1/2}
  \le\big((\velo^{\timelevel{n}})^{\transpose}\dualvolmat\,\velo^{\timelevel{n}}\big)^{1/2}
  +\big((\velo^{\stage{1}}-\velo^{\timelevel{n}})^{\transpose}\dualvolmat\,(\velo^{\stage{1}}-\velo^{\timelevel{n}})\big)^{1/2},
\end{equation}
because $\dualvolmat^{1/2}(\idmat-a\lapstag{k})^{-1}\dualvolmat^{-1/2}$
has norm at most one as well. The step therefore grows the energy by at most
the transported-to-old difference; moving walls may add their work through
$\velo^{\physbnd}_w$, and the split-step caveat in \cref{sec:viscosity:time}
still applies.

\emph{Walls.} For a boundary face $b$ of the cell $c=\cbnd{b}$ and the
unknowns $f$ of component $k$ on the sides of $c$ (side fractions
$\sidefrac{c}{f}$ of~\cref{eq:momentum:sides}), the wall connections are
\begin{equation}\label{eq:viscosity:walls}
  \wallcond{w}=\begin{cases}
    \dfrac{\meas{b}/2}{\abs{(\pos{f}-\pos{b})_{\faxis{b}}}},
      & \faxis{b}\neq k\text{ and the dual strip touches }b,\\[6pt]
    \dfrac{\sidefrac{c}{f}\meas{b}}{\cellwidth{c}{k}}, & \faxis{b}=k\ (\text{normal, the }f\text{ opposite }b),\\
    0, & \text{otherwise}.
  \end{cases}
\end{equation}
with the wall value $\velo^{\physbnd}_w$ at the point where the wall meets
the line of $f$: \code{NoSlip(u)} includes both kinds, with the value of
\code{u} (default zero) on tangential and zero on normal connections;
\apitype{FreeSlip} keeps the normal ones only (zero tangential stress); with
\apitype{VelocityBoundaryConditions} each side chooses, \apitype{Inflow}
with its function on both kinds, \apitype{Outflow} and
\apitype{ZeroGradientVelocity} with none, a homogeneous Neumann condition.
Connections of zero distance are skipped. Affine profiles, also between two
moving walls, are reproduced exactly, including on composite meshes: the
tendency of Couette flow between a stationary floor and a moving lid stays at
or below $2\times10^{-12}$ of its velocity maximum across a level transition
as well. A row next to a wall is not exact for quadratics: for $\velo_1=y(\domextent{2}-y)$ between no-slip walls it misses
$\Delta\velo_1=-2$ by $0.5$, independently of the mesh width (derivation
audit, \code{proof_review/eq-viscosity-walls/}).
The operator returned by \apifun{diffusion!} is
\begin{equation}\label{eq:viscosity:single-operator}
  \frac{\visc}{\dens}\Big((\lapstag{k}\velo)_f+\stagvol{f}^{-1}\sum_{w\in\wallconns{k},\ f(w)=f}\wallcond{w}\,\velo^{\physbnd}_w\Big),
  \qquad f\in\facesaxis{k}\cap\facesint,
\end{equation}
and zero on the boundary faces, whose normal values are prescribed.

\emph{Cut cells.} The unembedded and embedded operators share the geometric
contact construction. With a solid, the constructor of
\apitype{StaggeredDiffusionOperator} that takes a \apitype{GeometryField}
uses the cut-cell schedule of \cref{sec:momentum:dual} (weights times
apertures, open volumes $\opendualvol{f}$, one where they vanish) and rebuilds
the transverse connections from the open traces of the cells.
A normal connection between the two unknowns of one cell takes the open area
of that cell's mid-plane over the tangential overlap of the two faces: the
schedule average halves whenever one of the two primal faces flanking the
connection is closed, which is exactly the wall-adjacent row (\issue{110}).
The open trace of a cut cell is bisected at the cell centre, as on an uncut
cell, so a clipped trace does not move the boundary of a staggered dual
strip.
Across an open face $h$ of axis $t\neq k$ between two cells,
the extent of each cell along $k$ is clipped to the open part of its side on
$h$ and bisected, each half belonging to the unknown of that cell on its
side; every unknown of one cell that touches $h$ is connected to every such
unknown of the other cell, with the overlap of their halves within $h$ as
the measure and their distance along $t$ (on every mesh, also without level
transitions). The two segments of a split side of a cell, next to finer
cells, are connected along $t$ with half the open length of the line through
the split point. The derivatives across the connections are exact for linear
velocities: a fine--coarse connection along $k$ differences the coarse
unknown and the interpolation of the two fine siblings to its line; where
transverse connections are offset along $k$ (at a level transition), the
connections of an unknown to one side whose measure-weighted offsets along $k$
cancel share one centred derivative, with the coefficient $-1/d$ for the
unknown and $\meas{g}/(d\sum_{g'}\meas{g'})$ for each neighbour, $d$ the
measure-weighted offset along $t$; and an offset connection left without such
a group, where a level transition meets a wall normal to $k$, takes the zero
normal velocity of a \apitype{NoSlip} or \apitype{FreeSlip} wall as a
Dirichlet partner, interpolating the far unknown to the wall distance of the
near one (exact for linear profiles that vanish at the wall). An
\apitype{Inflow} side prescribes that velocity too, so it closes the connection
the same way with the prescribed value; because the wall values of the two
signed corrections sit at the two faces' own tangential positions, that form is
exact for inflow data constant along the wall and first order for data that
varies along it. At \apitype{Outflow} and \apitype{ZeroGradientVelocity} walls
no value is prescribed, so the connection is differenced on the tangential
level of one of its ends instead: the
endpoint that has a neighbour on its own level spanning the contact, and that
neighbour. Differencing the offset pair itself would carry the tangential
derivative of the offset into the $t$ derivative, with weight
$\Delta k/\Delta t$ and hence without vanishing under refinement; the
level-matched pair is exact for linear velocities (\issue{99}). The one-ring
face graph supplies the level-matched face; where none lies within two contact
distances, the weighted minimum-norm derivative of
\cref{eq:viscosity:gradient} over the two-ring face graph is used, as for the
unembedded transition. Both keep the coefficients summing to zero and the row
energy a non-negative square. The domain walls are~\cref{eq:viscosity:walls}
with each measure multiplied by the open fraction of the wall segment that
the connection owns (for a tangential connection the half of the cell's wall
segment on the side of $f$, owned whole by the segment of a split side that
touches the wall, for a normal one the part of the wall segment over $f$;
connections whose segment is covered are dropped). Embedded segments come
from cut cells and from uncut open cells sharing a face with a covered cell
(adjacent faces on one side are merged). For a Cartesian wall segment recovered
from the face aperture, a tangential face owns the part of the segment within
its dual strip; its
conductance is that length divided by the wall-normal distance to the face,
bounded below by one quarter of the smallest cell width. In particular, an
open tangential node a distance $h/2$ from a grid-aligned wall uses
$\meas{f}/(h/2)$, not a distance to the primal-cell segment midpoint.
A reconstructed PLIC segment of a cut cell, Cartesian or tilted, is split
among the dual faces of that cell by the tangential strip of a component face
and the half of the cell on that face's side, so the pieces partition the
segment and each piece is owned once. The wall distance of an unknown $f$ is
the \emph{perpendicular} distance $d_f$ to the reconstructed plane of its
cell, not the distance to the segment midpoint, and the conductance is
\begin{equation}\label{eq:viscosity:cut}
  \wallcond{w}=\frac{\meas{\interface^{\solid}_{c,f}}}
    {\max\roundpar{\abs{d_f},\ \tfrac{1}{64}\min_t\cellwidth{c}{t}}}
  \quad(\text{omitted for \apitype{FreeSlip}\code{()}}),
\end{equation}
with $\meas{\interface^{\solid}_{c,f}}$ the length of the piece. The floor is
one sixty-fourth of the smallest cell width and it is active: a cut plane
passes through or next to the centre of many wall-piece nodes, so a fifth to a
third of the pieces on the manufactured fixtures sit closer than $h/64$ (the
closest is a node exactly on the reconstructed plane). It bounds the largest
wall conductance at $64\meas{\interface^{\solid}_{c,f}}/\min_t\cellwidth{c}{t}$,
which keeps the assembled stiffness scaling as $1/h^2$; the pinned errors and
orders are measured with this floor in place. The factor
$\tfrac12\sidefrac{c}{f}$ of the previous closure is gone because the pieces
already partition the segment. A fluid-side unknown satisfies
$0=\velo_w=\velo_f-d_f\partial_n\velo$ for the wall-normal derivative
$\partial_n\velo$, so $\wallcond{w}\velo_f$ equals the wall flux
$\meas{\interface^{\solid}_{c,f}}\partial_n\velo$ of a linear profile
exactly. A node whose centre lies inside the solid gets the same positive
conductance, hence the opposite row contribution to the one its open dual
volume requires: it has no fluid on the wall side, and prescribing it would
remove a face that the cut-cell projection updates. Its rows therefore stay
first order with wall distance (\issue{110}); the open-midpoint LS-STAG
closure of \cref{sec:viscosity:lsstag} removes this defect.
Where a grid-aligned wall coincides with a normal component face, that face is
prescribed zero and removed from the diffusion unknowns. The mid-plane contact
of the adjacent cell is fully open there and carries the whole Dirichlet flux
to the eliminated face, for both wall types; no separate wall connection is
added, because an aperture-weighted half contact plus a compensating half
connection is the same matrix split differently. The retained principal
stiffness block also
carries flux to the eliminated face. Its integrated reaction, the negative
column sum of the eliminated-to-retained stiffness block applied to the
retained velocities, is stored separately from wall connections and included
when reporting embedded viscous forces; it is not added to diffusion twice.

Every embedded wall value is zero, also for a moving \code{NoSlip(u)}
(since \issue{75}). With \apitype{VelocityBoundaryConditions} the embedded
walls take its separate \code{embedded} condition, whatever the sides choose:
\apitype{NoSlip}\code{()} by default, or \apitype{FreeSlip}\code{()}, which
omits their tangential connections as uniform free slip does. Away from
eliminated wall-normal faces, a completely covered background face retains the
unit-volume algebraic penalty $\wallcond{w}=1$ for both wall types. With
$\velo^{\timelevel{n}}_f=0$, the midpoint step
\cref{eq:viscosity:single-midpoint} scales such a covered value by
$1/(1+\tfrac12\dt\visc/\dens)$; the flow solver also sets closed-face values
to zero at the start of every step and after projection
(\apifun{enforce_geometry_no_slip!}). Solid surfaces on cell faces are included
as well as cut walls (\issue{79}). The aligned-wall manufactured Helmholtz
solution now converges at second order on uniform and composite meshes in the
maximum, $L^2$ and $L^1$ norms (\issue{106};
\code{test/operators/embedded_wall_rows.jl}). A cut-cell PLIC segment now uses
its own reconstructed plane, and the wall flux of a fluid-side unknown is
exact for a linear velocity, but the manufactured Helmholtz solution on a
tilted wall still converges with observed $L^2$ orders of $1.1$ to $1.4$, and
$0.74$ on the composite grid, rather than two (\issue{110};
\code{test/operators/embedded_wall_rows.jl}). The remaining first-order error
sits in the cut cells' interior contact measures, which are rebuilt from
bisected one-dimensional traces: the normal component's wall rows keep an
$h/(2\abs{\nabla d})$ residual, $0.091$ rms at $32$ cells on the two slopes; a node whose centre lies inside the solid
also carries its wall piece with the flux sign of a fluid-side node, which a
symmetric non-negative stiffness cannot avoid.
The constructor builds when a solid cuts cells at a level transition
(\issue{78}). Rows at a level transition can have a truncation error of order
one; the derivation audit measured $0.09$ for $\abs{\Delta\velo}\le0.31$ on
its composite mesh, unchanged from $16$ to $64$ cells across while the other
rows converged with order two. The solution of a Helmholtz problem with level
transitions meeting no-slip or free-slip walls nevertheless converges with
order at least $1.8$ in the maximum norm (\issue{87}).

\begin{sloppypar}
\emph{Checks} (\code{figures/viscosity.jl}). On a uniform periodic mesh an
independent five-point stencil equals \apifun{diffusion!} to
$\chkViscUniform$ (relative; \code{uniform_laplacian_check}). On the walled
composite mesh of \cref{ch:forces}, the weighted asymmetry is
$\chkViscWeightedAsym$ and the smallest stiffness eigenvalue relative to the
largest entry is $\chkViscMinEig$ (\code{symmetric_laplacian_check}).
On periodic, channel and walled composite fixtures and
$\dt\visc/\dens\in\{10^{-3},10^{-1},10,10^3\}$, the largest backward-Euler
and midpoint energy factors were $1.663$ and $4.408$ before the change,
and $1+5.7\times10^{-13}$ and $1+4.5\times10^{-13}$ after it.
\end{sloppypar}

\begin{deviation}{Single-phase viscous operator}\label{dev:viscosity:single}
The thesis discretises the stress $\divergence(2\visc\strainrate)$ with
staggered strain operators on conforming meshes~\thesis{eq.~(7.52)} and
defers refinement interfaces to Van der Plas~\cite{vanderplas2017}. For a
single phase the code uses the componentwise Laplacian $(\visc/\dens)\Delta$,
equal to the stress form for constant viscosity and $\divergence\Velo=0$,
with shared geometric contact energies at level transitions. These are the
package's own stencils.
\end{deviation}

\begin{implnote}
\apitype{StaggeredDiffusionOperator}\code{(mesh; viscosity=1, boundary=NoSlip())}
(\srcsym{src/operators/staggered_diffusion.jl}{StaggeredDiffusionOperator})
and \code{(geometry; viscosity, boundary)}
(\srcsym{src/embedded/operators.jl}{StaggeredDiffusionOperator}); the flow
solver passes the kinematic viscosity $\visc/\dens$
(\srcsym{src/navier_stokes/system.jl}{NavierStokesSystem}).
Connections \srcsym{src/embedded/operators.jl}{_geometric_diffusion_connections}~\cref{eq:viscosity:connections},
derivatives \srcsym{src/embedded/operators.jl}{_cut_cell_gradient_stencil}~\cref{eq:viscosity:gradient},
stiffness \srcsym{src/embedded/operators.jl}{_geometric_diffusion_stiffness}~\cref{eq:viscosity:stiffness},
and $\lapstag{k}=-\dualvolmat^{-1}\viscstiff{k}$~\cref{eq:viscosity:laplacian},
walls \srcsym{src/operators/staggered_diffusion.jl}{_wall_diffusion_connections},
\srcsym{src/operators/staggered_diffusion.jl}{_include_wall_connection} and
\srcsym{src/operators/staggered_diffusion.jl}{_wall_value}~\cref{eq:viscosity:walls},
cut cells \srcsym{src/embedded/operators.jl}{_cut_cell_diffusion_component}
and \srcsym{src/embedded/operators.jl}{_embedded_wall_connections}~\cref{eq:viscosity:cut}
with the segments on cell faces \srcsym{src/embedded/operators.jl}{_embedded_aligned_wall_segments},
transverse connections \srcsym{src/embedded/operators.jl}{_geometric_diffusion_connections}
with \srcsym{src/embedded/operators.jl}{_open_dual_trace} and
\srcsym{src/embedded/operators.jl}{_add_split_side_diffusion_connections!},
derivatives \srcsym{src/embedded/operators.jl}{_cut_cell_gradient_stencil},
\srcsym{src/embedded/operators.jl}{_cut_cell_transverse_stencils},
\srcsym{src/embedded/operators.jl}{_wall_contact_stencil!} and
\srcsym{src/operators/staggered_diffusion.jl}{_level_matched_gradient} with
\srcsym{src/operators/staggered_diffusion.jl}{_minimum_norm_gradient},
the clipped domain-wall measure \srcsym{src/embedded/operators.jl}{_wall_connection_measure}.
\apifun{stable_timestep}\code{(operator; cfl=0.9)} returns the explicit bound
$\code{cfl}\,\dens/(\visc\max_f\sum_{f'}\abs{(\lapstag{k})_{f,f'}})$, which
the implicit flow step does not use. Tests:
\testsym{test/operators/staggered_diffusion.jl}{staggered diffusion stable core on uniform and AMR meshes},
\testsym{test/operators/staggered_diffusion.jl}{composite diffusion solution convergence (#80)},
\testsym{test/operators/staggered_diffusion.jl}{composite implicit diffusion is energy contractive (#80)},
\testsym{test/operators/staggered_diffusion.jl}{uniform periodic staggered Laplacian convergence},
\testsym{test/operators/staggered_diffusion.jl}{wall slip and no-slip closures},
\testsym{test/indicator/geometry.jl}{Embedded free slip omits only physical wall-shear stresses},
\testsym{test/operators/staggered_diffusion.jl}{moving domain walls do not move an embedded disc}
and \testsym{test/operators/staggered_diffusion.jl}{domain-wall diffusion uses only uncovered wall segments}
(\issue{75}); cut cells
\testsym{test/operators/staggered_diffusion.jl}{cut-cell diffusion connections at an AMR transition (#78)},
\testsym{test/operators/staggered_diffusion.jl}{cut transverse measures use overlapping traces (#87)},
\testsym{test/operators/staggered_diffusion.jl}{wall conductance follows split dual strips (#87)},
\testsym{test/operators/staggered_diffusion.jl}{wall-contact flux respects a linear normal profile (#87)},
\testsym{test/operators/staggered_diffusion.jl}{cut-cell Helmholtz convergence at walls and level transitions (#87)}
(order at least $1.8$ in the maximum norm),
\testsym{test/operators/staggered_diffusion.jl}{orphan transverse rows are linear-exact at open walls (#99)},
\testsym{test/operators/staggered_diffusion.jl}{open sides share the orphan closure (#99)},
\testsym{test/operators/staggered_diffusion.jl}{inflow wall-contact closure is exact for constant wall data (#99)}
and \testsym{test/operators/staggered_diffusion.jl}{inflow wall-contact closure is first order for varying wall data (#99)}
and \testsym{test/operators/staggered_diffusion.jl}{orphan transverse rows converge on a manufactured field (#99)},
\testsym{test/two_phase/variable_staggered_diffusion.jl}{Grid-aligned embedded walls have the cut-wall limit}
(\issue{79}, also for the operator of \cref{sec:viscosity:stress}).
\end{implnote}

\subsection{LS-STAG closure from PLIC data}\label{sec:viscosity:lsstag}
The LS-STAG closure is the default: with the keywords omitted,
\code{StaggeredDiffusionOperator(...)}
and \code{NavierStokesOptions()} select \code{closure=LSStagClosure()} and
\code{embedded_viscous_closure=LSStagClosure()}. The rows of the faces next to
cut cells are then replaced by the LS-STAG
closure of Cheny and Botella~\cite{cheny2010} (\issue{110}, \issue{13}),
assembled from the PLIC data of the \apitype{GeometryField}. The explicit
\code{WallDistanceClosure()} selects the legacy closure described above.

The default was switched to LS-STAG after measuring both closures, and the
owner judged the regressions benign. Developed $\mathrm{Re}=100$ flow around a
cylinder at 16 cells per diameter gives mean drag, pressure drag, maximum lift
and Strouhal number $1.25876$, $0.92661$, $0.2154$, $0.1447$ with LS-STAG
against $1.14227$, $0.83327$, $0.1534$, $0.1495$ with wall distance
(reference $1.24$, $0.93$, $0.30$, $0.165$): pressure drag and lift move
towards the reference, viscous drag and shedding frequency do not. A
zero-gravity sessile cap on a $15^\circ$ slope at a quarter capillary time
changes its open-volume-fraction shape error by $+6.70\,\%$ at $60^\circ$ with
16 cells across the domain and $-13.73\,\%$ at 32 cells, and by $+0.21\,\%$
and $+1.54\,\%$ at $90^\circ$, where the kinetic energy grows by $2.36\,\%$
and $3.52\,\%$.

For a single-phase component $k$ the \emph{replaced faces} are the faces of the cut cells (solid fraction
strictly between zero and one), their neighbours on the same grid line (the
line $x_k=\text{const}$ through the face) and, recursively, the fine faces
whose grid line ends at a coarse cell with a replaced face; a cell with a
replaced face is a \emph{Green cell}. Every connection with a replaced end,
every normal connection inside a Green cell, and the wall and boundary
connections of replaced faces are removed and replaced by the terms below.
The unknowns, the mass $\dualvolmat$ (the open staggered volumes), the
divergence and the gradient are unchanged, and on uncut cells of one level
both closures give the same two-point terms, so the seam between them is
consistent.

\emph{Values.} A face with an open staggered volume and an aperture
$a_f>10^{-10}$ is open. Its open part has the length $\ell_f=a_f\meas{f}$ and
lies at the end that the solid planes of its cut cells leave open: the first
cut cell whose own trace of the face is partial decides by the sign of the
tangential component of its normal. The value $\velo_f$ is the velocity at
the midpoint $\mathbf m_f$ of the open part, which lies in the fluid also
when the face centre does not.

\emph{Normal stress.} A Green cell $c$ with open area $\abs{c}$ has the Green
gradient and energy
\begin{equation}\label{eq:viscosity:lsstag}
  G_c\velo=\frac{1}{\abs{c}}\sum_{f\in\partial c}o_{c,f}\,\ell_f\,\velo_f ,
  \qquad
  E^{\mathrm n}_c=\abs{c}\,\roundpar{G_c\velo}^2 ,
\end{equation}
with $o_{c,f}=\pm1$ the side of $f$ and the wall value zero; $G_c$ is the
open divergence applied to one component. A side of $c$ on the domain
boundary enters with its prescribed normal velocity as a source; an outflow or
zero-gradient side leaves the cell without a normal term, as before. The wall
closes the sum: its reaction $o_{c,f}\ell_f\sum_{f'}o_{c,f'}\ell_{f'}/\abs{c}$
per unit $\velo_f$ is stored with the wall reactions and enters the viscous
force on the solid.

\emph{Shear stress.} Along the grid line of $f$, the half control volume of
$f$ in each adjacent cell owns half of the open length of the lower and upper
edge of its strip in that cell, $T^{\mathrm{lo}}$ and $T^{\mathrm{hi}}$ (the
bisection by measure that also defines the open volumes). An open end of $f$
that meets an open end of the next face $f'$ on the line is a link with
area $A=\sum_{\text{sides}}\min(T^{\mathrm{hi}}_f,T^{\mathrm{lo}}_{f'})$ and
energy $A\,d\,((\velo_{f'}-\velo_f)/d)^2$, $d$ the distance of the two
midpoints. A line that ends at a coarse cell takes the derivative along the
line that is exact for affine functions through $\mathbf m_f$ and the faces of
the coarse cell on the two neighbouring lines (a closed face contributes a
point of the wall with value zero, a domain boundary its prescribed normal
velocity) with the energy $A\abs{w_f}^{-1}(D\velo)^2$, $w_f$ the weight of
$\velo_f$, so the partners receive their share of the flux. The wall, or the
domain boundary at an open end on it, carries what the links leave of the
projected extent $\max(T^{\mathrm{lo}},T^{\mathrm{hi}})$ of each half control
volume, by the one-sided quotient to the wall value at $\delta=\ell_f/2$ at a
closed end and at the distance along the line to the PLIC plane beyond an open
end. Every half control volume is then closed, $\oint n_t\,\mathrm dS=0$, and
an affine velocity that vanishes on the wall has zero residual (largest
relative residual $10^{-10}$ in the tests). As in~\cref{eq:viscosity:cut} the
embedded quotients are omitted for \code{FreeSlip()}. A face of a cut cell
without an open part but with an open volume (next to a covered cell) keeps the
wall value through the diagonal $\stagvol{f}(2/h_k^2+2/h_t^2)$.

Every term is a non-negative weighted square, so $\viscstiff{k}$ stays
symmetric positive semidefinite, and definite with a wall; no face is removed
and no distance floor is used: the open midpoint lies in the fluid, and its
distance along the line to a closed end is exactly $\ell_f/2$. On the \#110
Helmholtz fixture of \cref{eq:viscosity:cut} (wall $y=1+0.1x$, $n=16$ to
$128$, errors at the open midpoints over every open face) the largest error
of component~1 at $n=128$ drops from $2.2\times10^{-2}$ (orders
$0.97/0.91/0.99$) to $1.8\times10^{-5}$ (orders $2.00/2.00/2.00$). On the
slopes $0.1$ and $0.5$ and on a circle, uniform and with the right half
refined, the $L^2$ and $L^1$ orders are $1.91$ to $2.07$ and the maximum-norm
orders $1.68$ to $2.00$, except in the corner where the slope-$0.5$ wall
meets the boundary $x=8$ under an interior angle of $116.6^\circ$
($1.31/1.54/1.77$, $\pi/\omega=1.54$).

\emph{Small cells.} Without a floor a thin layer has the conductance of its
width: $h^2$ times the largest absolute row sum of $\lapstag{k}$ reaches
$7.8\times10^{5}$ for an open cell fraction of $3\times10^{-8}$, against
$3.1\times10^{4}$ with the floor of~\cref{eq:viscosity:cut}. The flow step is
implicit, so this bounds no time step. The damped-Jacobi weight of the
viscous V-cycle depends on the ratio of a row's absolute sum to its diagonal,
not on its size, and stays at $0.66$ to $0.73$ on all fixtures (the wall-distance
closure drops to $0.20$ where its own floored rows are small). Over
$n=32$ to $256$, uniform and composite, diffusion numbers $\dt\visc/(2h^2)$
from $0.16$ to $10^3$, a cut circle and a slope and the offsets of each that
leave the smallest open cell within one cell width (open fractions down to
$1.6\times10^{-8}$), the preconditioned CG iterations with the LS-STAG closure
are $0.15$ to $1.74$ times those of the wall-distance closure: at most $1.74$ on the
smallest cells (circle, composite, $n=64$: $54$ against $31$) and as low as
$0.15$ where the wall-distance smoother weight drops ($15$ against $99$). A floor
would bound these rows at the price of an $O(1)$ relative flux error in
exactly the cells it touches.

\emph{Scope.} These componentwise rows apply to single-phase diffusion.
Variable viscosity uses the coupled rows below. The solid itself is
stationary: a moving domain wall does not prescribe motion to an embedded
wall. With a free-slip solid the single-phase componentwise operator is not
consistent at an oblique wall for either closure. The wall reaction of a
single-phase Green cell on an inflow side omits the prescribed inflow term.

\begin{implnote}
\code{LSStagClosure} and \code{WallDistanceClosure}
(\srcsym{src/operators/embedded_viscous_closure.jl}{LSStagClosure});
assembly \srcsym{src/embedded/ls_stag.jl}{_ls_stag_closure}~\cref{eq:viscosity:lsstag}
with the open parts \srcsym{src/embedded/ls_stag.jl}{_ls_stag_open_part} and
the owned strip edges \srcsym{src/embedded/ls_stag.jl}{_ls_stag_strip_edges},
called from \srcsym{src/embedded/operators.jl}{_cut_cell_diffusion_component}.
Tests:
\testsym{test/operators/embedded_ls_stag.jl}{LS-STAG closure is second order at oblique walls (#110)},
\testsym{test/operators/embedded_ls_stag.jl}{LS-STAG closure keeps the unknowns, the mass and an SPD stiffness (#110)},
\testsym{test/operators/embedded_ls_stag.jl}{LS-STAG rows of an affine field that vanishes on the wall are exact (#110)},
\testsym{test/operators/embedded_ls_stag.jl}{LS-STAG viscous multigrid versus wall distance (#110)}
and \testsym{test/operators/embedded_wall_force.jl}{cut-cell wall force balances viscous momentum loss for both closures (#110)}.
\end{implnote}

\section{Two-phase stress operator}\label{sec:viscosity:stress}
The two-phase operator discretises $\divergence(2\visc\strainrate)$ as the
negative adjoint of a strain-rate map. Its viscosity is, in every cell, a
mixture of the phase viscosities with the conditional fraction
$\condfrac{\liq}{c}$~\thesis{\S7.5.1}. The default is the harmonic (series)
mixture and the thesis' geometric mixture is the
\apitype{GeometricViscosityAveraging} option,
\begin{equation}\label{eq:viscosity:viscosity}
  \visc_c=\begin{cases}
    \roundpar{\dfrac{\condfrac{\liq}{c}}{\visc^{\liq}}+\dfrac{1-\condfrac{\liq}{c}}{\visc^{\gas}}}^{-1}
      & \text{(harmonic, default)},\\[8pt]
    (\visc^{\liq})^{\condfrac{\liq}{c}}(\visc^{\gas})^{1-\condfrac{\liq}{c}}
    =\exp\roundpar{\condfrac{\liq}{c}\log\visc^{\liq}+(1-\condfrac{\liq}{c})\log\visc^{\gas}}
      & \text{(geometric, \apitype{GeometricViscosityAveraging})}.
  \end{cases}
\end{equation}
The rows of the strain map $\strainmat:\facespace\to\mathbb{R}^{\text{rows}}$
are of three kinds:
\begin{equation}\label{eq:viscosity:rows}
  \begin{aligned}
  &\text{(normal)} && (\strainmat\velo)_{c,k}=\frac{1}{\meas{\opencell{c}}}\sum_{f\in\faceset{c}\cap\facesaxis{k}\cap\facesint}\orient{c}{f}\,\aperture{f}\meas{f}\,\velo_f,
    && \openfrac{c}>0,\ k=1,2,\\
  &\text{(shear)} && (\strainmat\velo)_r=(\vderiv{2}\velo_1)_{g}+(\vderiv{1}\velo_2)_{g'},
    && r\in\shearrows,\ \shearregion{r}=\shearregion{g}\cap\shearregion{g'},\ \meas{\shearregion{r}}^{\open}>0,\\
  &\text{(wall)} && (\strainmat\velo)_w=\velo_{f(w)},
    && w\text{ wall or covered unknown},
  \end{aligned}
\end{equation}
where $\velo_1,\velo_2$ are the two components. A normal row is the
directional open divergence $\partial_k\velo_k$ of the cell; prescribed
normal wall faces have zero velocity. At a Cartesian embedded wall their
columns, and the adjacent covered-face columns, are eliminated from the
physical strain map. Their global face slots remain isolated unit algebraic
rows, with zero right-hand side and solution, so the coupled operator keeps
the face-field ordering. The eliminated strain columns are retained
separately for wall-reaction evaluation as viscosity changes. Wall rows exist
for cut cells and uncut open cells sharing a face with a covered cell
(\issue{79}), and for covered unknowns outside the constrained neighbourhood.
Cartesian embedded-wall shear is represented by tangential wall rows; the
corresponding cell quarters are omitted from the shear-row measures to avoid
counting shear across the solid twice. The normal cell stress already supplies
the normal flux. For free slip the tangential wall rows are omitted.

\emph{Shear partition.} For component $k$ and the tangent $t\neq k$, the
samples are the unknowns $f\in\facesaxis{k}\cap\facesint$ and the wall faces
normal to $\unitvec{k}$, whose velocity is zero. The staggered control volume
of a sample spans the face along $t$ and half of each adjacent cell along
$k$; these volumes tile the domain. Where the upper side (along $t$) of the
volume of a sample $a$ meets the lower side of the volume of a sample $b$,
their overlap is a contact $g=(a,b)\in\shearcontacts{k}$ of measure
$\meas{g}$; a sample whose volume ends on a wall normal to $\unitvec{t}$ has a
wall contact with that wall, of measure the width of the volume. With
$\conndist{g}$ the distance along $t$ between the centre of $a$ and the centre
of $b$ (or the wall), the region of the contact is the rectangle
\begin{equation}\label{eq:viscosity:shear-partition}
  \shearregion{g}=(\text{overlap along }k)\times(\text{centre of }a,\ \text{centre of }b)_t,
  \qquad
  \meas{\shearregion{g}}=\meas{g}\conndist{g},
  \qquad
  \textstyle\bigcup_{g\in\shearcontacts{k}}\shearregion{g}=\domain ,
\end{equation}
a partition of the domain (periodic images identified), because a region is
the upper half of the volume of $a$ and the lower half of the volume of $b$
over their overlap. The code works on the doubled finest lattice, on which
every face centre is an integer point, so the contacts and the intersections
below are exact. The tangential derivative of a contact is
\begin{equation}\label{eq:viscosity:shear}
  (\vderiv{t}\velo_k)_g=\begin{cases}
    0, & g\text{ a wall contact},\ \apitype{FreeSlip},\\[2pt]
    \mp\velo_a/\conndist{g}, & g\text{ a wall contact},\ \apitype{NoSlip}\ (\text{upper, lower wall}),\\[2pt]
    (\velo_b-\velo_a)/\conndist{g}, & a,b\text{ paired},\\[4pt]
    \pm\Big(\dfrac{1}{M_a}\displaystyle\sum_{g'=(a,b')}\meas{g'}\velo_{b'}-\velo_a\Big)\Big/\bar d_a,
      & \text{anchored at }a,
  \end{cases}
\end{equation}
where wall samples enter with the value zero and a wall contact takes the
condition of the side it meets (each side its own with
\apitype{VelocityBoundaryConditions}). Two samples are paired if
their volumes have equal widths along $k$ (a wall sample counted with its
mirror image) or if the narrower one has no other contact on that side; the
code requires paired samples to be collinear. Otherwise, at a level
transition across $\unitvec{t}$, the narrower sample $a$ (a fine segment
between two coarse unknowns) anchors all its contacts $g'=(a,b')$ on that
side, with $M_a=\sum_{g'}\meas{g'}$,
$\bar d_a=\sum_{g'}\meas{g'}\conndist{g'}/M_a$ and the sign $+$ when the
partners lie on the upper side; the code requires the partners to be wider
and centred, $\sum_{g'}\meas{g'}(\pos{b'}-\pos{a})_k=0$. Each derivative is
then exact for affine velocities (rows with a wall value for fields that
vanish there), and every derivative of a contact without a wall sample
annihilates constants. These are the fine--coarse closures of the
single-phase operator, \cref{eq:viscosity:gradient}.

\emph{Shear rows.} A shear row $r\in\shearrows$ is a nonempty intersection
$\shearregion{r}=\shearregion{g}\cap\shearregion{g'}$ of a contact $g$ of
$\velo_1$ and a contact $g'$ of $\velo_2$, whose strain rate is the sum of the
two derivatives~\cref{eq:viscosity:rows}; rows whose two derivatives both
vanish identically are omitted. The rows partition the domain once more: they
are the common refinement of the two contact partitions. On a uniform mesh
both partitions are the cells of the vertex-centred dual mesh, and every row
is the classical vertex row with a two-point difference for each component.
With the measures of \cref{eq:viscosity:weights} (uncut) the force of the
shear rows on an unknown $f$ of component $k$ is a sum of contact fluxes,
\begin{equation}\label{eq:viscosity:shear-flux}
  -\sum_{r\in\shearrows}\meas{\shearregion{r}}\visc_r(\strainmat\velo)_r\,\strainmat_{r,f}
  =\sum_{g\ni f}\pm\,\ell_{g,f}\,\bar\tau_g,
  \qquad
  \bar\tau_g=\frac{1}{\meas{\shearregion{g}}}\sum_{\shearregion{r}\subset\shearregion{g}}\meas{\shearregion{r}}\visc_r(\strainmat\velo)_r ,
\end{equation}
because the rows inside a contact region add up to
$\meas{\shearregion{g}}=\meas{g}\conndist{g}$, which cancels the factor
$1/\conndist{g}$ of the derivative: the length is $\ell_{g,f}=\meas{g}$ for a
paired or wall contact. The contacts $g'$ of an anchored group share one
derivative, so for them $\bar\tau_g$ is the stress averaged over the whole
group, $\bar\tau_{a}=\sum_{g'}\meas{\shearregion{g'}}\bar\tau_{g'}/(M_a\bar d_a)$:
the anchor receives $M_a\bar\tau_a$ and each partner $b'$ receives
$\meas{g'}\bar\tau_a$, not the average over its own region. The group regions
lie within $O(h)$ of the anchor, so $\bar\tau_a$ is a first-order consistent
stress at every partner's contact as well. Each staggered volume therefore
receives on its two tangential sides region- or group-averaged shear stresses
times lengths that add up to its width. For a constant viscosity and an
affine velocity every average is the exact stress, so the shear forces cancel
on every volume whose contacts are free of wall values, at level transitions
included; for
smooth fields the fluxes are first-order accurate at transitions and second
order elsewhere, which gives a second-order solution (checks below). The
identity needs the measure of the rows of a contact region to be
$\meas{g}\conndist{g}$ for both components at once. The quarter-cell vertex
rows used before \issue{112} violate it at a transition: on a straight
transition with coarse width $2h$ the contacts of both components through a
coarse corner and through the hanging vertex between two corners have the
regions $\tfrac32h\times h$, while the quarter cells give
$\tfrac52h^2$ and $\tfrac12h^2$; the shear forces of an affine velocity did
not cancel, the truncation error grew as $h^{-1}$ along the transition and
the solution converged at first order. At the corners of a refined region
the regions of the two components differ, so no measure per vertex is
consistent for both, and the rows are their intersections.

\paragraph{Two-phase LS-STAG at a stationary solid.}
The normal rows are the open Green gradients in \cref{eq:viscosity:rows}.
At an open vertex $v$, let $D_v u$ denote the \emph{sum}
$\partial_y u_x+\partial_x u_y$, formed from the two open-midpoint
differences. Where a fine grid line ends, the directional derivative uses
the affine stencil through the neighbouring coarse faces if both are open.
When a coarse partner is covered, or the fluid line ends exactly at a grid
vertex on the solid, both boundary conditions instead take the endpoint
trace from the fluid side; no shared shear rectangle straddles a missing
sample. An endpoint clipped in the interior of a face has zero prescribed
value in the uncondensed Green map.

For an open cell of area $V_c$, define $Q_c$ as the outward flux of
$(u_y,u_x)$ through its boundary. Each Cartesian edge uses its open length
times the average of the two endpoint traces. Its two adjacent cells
therefore receive opposite fluxes, including at a split side. A complete
regular cell has four vertex contributions of area $V_c/4$. Elsewhere the
raw corner area is the product of the incident open half-edge lengths; a
hanging or closed corner without both lengths has no such rectangle.
Scale the areas together, if necessary, to leave at least $V_c/2$ for the
remaining part of a cut or transition cell. With the resulting areas $a_{cv}$
and $R_c=V_c-\sum_v a_{cv}>0$, its additional shear row is
\begin{equation}
    D_c^{\mathrm{rem}}u =
    \frac{Q_c(u)-\sum_v a_{cv}D_vu}{R_c}.
    \label{eq:viscosity:lsstress}
\end{equation}
Thus the area-weighted strains sum to the cell Green flux for each choice
of boundary trace. The $a_{cv}$ also give the viscosity interpolation weights
of the shared vertex; the remaining row uses the cell viscosity.
For an uncut uniform cell $R_c=0$ and the
four vertex rows already give the Green flux. In the Cartesian free-slip
limit there is no wall-shear row. If a slip wall ends at a grid vertex,
an adjacent full cell uses the one-sided fluid trace extrapolated from open
midpoints; setting it to zero would damp even a constant translation
tangent to the wall. At an oblique free-slip PLIC segment a scalar
tangential trace $s$ is free, with normal velocity fixed to zero. If $E_c$
collects the two normal rows and the remaining shear row before elimination,
their strain is $E_c u-b_c s$. On a planar wall admitting a rigid unit
tangent $u_t$, the wall column satisfies $b_c=E_c u_t$: the shared
derivatives $D_v$ annihilate constants, and $b_c$ uses the same endpoint
traces as $Q_c$. The local minimum of the three weighted strain squares over
$s$ therefore preserves this motion and the positive factorization, without
a wall-distance penalty. The uncondensed identity is checked in every cut
cell of uniform and composite periodic planar bands in
\testsym{test/two_phase/variable_staggered_diffusion.jl}{Uncondensed LS-STAG FreeSlip rows reproduce the PLIC tangent}.

At each Cartesian box-wall vertex $v$ use one tangential velocity $s_v$ in
\emph{both} terms of \cref{eq:viscosity:lsstress}: the endpoint trace in
$Q_c$ is $s_v$, and the wall derivative in $D_v$ is
$\pm(u_f-s_v)/d$, where $u_f$ is the closest open midpoint at distance $d$.
The normal component is prescribed; both components are prescribed at a box
corner. Stationary no slip sets $s_v=0$, including in cells where a level
transition meets the box. Free slip minimizes the weighted strain work over
the $s_v$. A trace seen only by its vertex releases that square. If several
traces enter one completion, their minimization is coupled; the cell and
vertex viscosities need not agree, so the minimum is updated when they change.

To see the stiffness ordering, write the complete no-slip energy as
$\|W^{1/2}Eu\|_2^2$. Let $B$ contain the box and tilted PLIC trace
columns. For each aligned embedded wall whose remaining shear row is
omitted in free slip, also include the coordinate vector of that row:
minimizing over its coefficient deletes exactly that square. If $P$ is the
orthogonal projector onto the range of $W^{1/2}B$, the minimized energy is
$\|(I-P)W^{1/2}Eu\|_2^2$. Hence the released energy is
$\|PW^{1/2}Eu\|_2^2\ge0$ for every $u$, even when the trace columns are
dependent. The implementation first eliminates the cell-local tilted
PLIC trace from \emph{both} $E$ and the box columns, then the remaining
shared box traces in their connected groups
(\srcsym{src/two_phase/diffusion/variable_diffusion_box_traces.jl}{_ls_update_box_strain!}).
A direct weighted projection,
including the aligned-wall case and refreshed viscosities, is checked in
\testsym{test/indicator/geometry.jl}{Embedded free slip omits only physical wall-shear stresses};
\testsym{test/two_phase/variable_staggered_diffusion.jl}{LS-STAG free-slip energy nests at terminating grid lines}
checks the stiffness ordering on uniform and composite planar bands.
On the right-half-refined Helmholtz fixture,
\testsym{test/two_phase/variable_staggered_diffusion.jl}{LS-STAG NoSlip Helmholtz convergence at box-boundary T-junction},
the no-slip maximum error at $n=32,64,128$ is
$9.76\times10^{-5},\,2.65\times10^{-5},\,6.99\times10^{-6}$
(orders $1.88,1.93$). Using a wall midpoint quotient only at the vertex
and a fluid extrapolation in the adjacent Green remainder instead gave
$6.57\times10^{-5}$ at $n=128$ and order $1.03$.

\paragraph{Per-side walls and corners.}\label{sec:viscosity:per-side}
With \apitype{VelocityBoundaryConditions} each box side is a stationary
\apitype{NoSlip}\code{()} or \apitype{FreeSlip}\code{()}, and the solid has
its own \code{embedded} condition, a stationary no-slip wall by default, as in
the single-phase operator. A no-slip solid keeps the rows of uniform no slip;
a free-slip solid releases its tilted trace and omits its aligned-wall rows
as uniform free slip does.
The trace $s_v$ of a box vertex $v$ that is not a corner is the velocity
component tangential to the one side that contains $v$, and that side alone
fixes $s_v=0$ or releases it
(\srcsym{src/two_phase/diffusion/variable_diffusion_ls_stag.jl}{_ls_box_entry!});
the legacy rows apply the same rule to the wall contacts
of~\cref{eq:viscosity:shear}.

\emph{Corner rule.} A box corner lies on two sides, and each velocity
component there is the normal velocity of one of them: $\velo_1$ of the side
normal to $\unitvec{1}$, $\velo_2$ of the other. No-slip and free-slip sides
both prescribe their normal velocity, zero, so the tangential trace of one
side at a corner is the prescribed normal velocity of the other, and the
corner velocity vanishes for every combination of the two. The corner line
takes its endpoint from the boundary face, whose column is not an unknown,
so no trace column exists there. Releasing a corner trace with its
free-slip side would let that side move fluid through the other side.

\emph{Nesting.} For LS-STAG the rows, their measures and their viscosities
do not depend on the side conditions; only the admissible traces do. Let
$\mathcal S_{\mathrm N}=\{0\}$, $\mathcal S$ and $\mathcal S_{\mathrm F}$ be
the box-trace spaces of all no-slip sides, of a mixed configuration and of
all free-slip sides, each with the corners fixed, so
$\mathcal S_{\mathrm N}\subseteq\mathcal S\subseteq\mathcal S_{\mathrm F}$,
and let $\mathcal P$ be the solid's parameters: $\{0\}$ for a no-slip solid,
the tilted traces and aligned-row coefficients of the projection argument
above, unconstrained, for a free-slip one, with columns $C$. With $E$, $W$
and the box columns $B$ of the previous paragraphs, every configuration is
the minimum over its spaces, and a minimum over larger sets is smaller:
\begin{equation}\label{eq:viscosity:per-side-nesting}
  \begin{gathered}
  \velo^{\transpose}K_{\mathcal S,\mathcal P}\velo
  =\min_{s\in\mathcal S,\,p\in\mathcal P}\|W^{1/2}(E\velo-Bs-Cp)\|_2^2,\\
  \mathcal S\subseteq\mathcal S',\ \mathcal P\subseteq\mathcal P'
  \ \Longrightarrow\ K_{\mathcal S,\mathcal P}\succeq K_{\mathcal S',\mathcal P'}.
  \end{gathered}
\end{equation}
All no-slip walls are $(\mathcal S_{\mathrm N},\{0\})$ and uniform
\apitype{FreeSlip}\code{()} is $(\mathcal S_{\mathrm F},\mathcal P)$, so every
per-side configuration, with either solid, lies between them, and releasing
the solid at fixed sides only removes dissipation. The shear rows
of~\cref{eq:viscosity:shear} (legacy closure, or no solid) drop the one-sided
difference of a free-slip wall contact instead and nest in the same
measurements. Over all sixteen side combinations, each with a no-slip and
with a free-slip solid, on fixtures without a solid, with the sloping, an
aligned and a disk solid, a $1.5h$ gap and a disk that cuts a corner cell, on
uniform and right- or top-refined meshes, with both closures and viscosity
ratios up to $1000$, the smallest eigenvalue of every difference along these
chains, and of every stiffness, relative to the largest entry of the
all-no-slip stiffness is $-3.8\times10^{-16}$
(\testsym{test/two_phase/variable_staggered_diffusion.jl}{Mixed per-side walls nest between no-slip and free-slip walls},
\testsym{test/two_phase/variable_staggered_diffusion.jl}{Per-side walls with a free-slip solid nest and keep its rigid tangent}).
Equal sides and solid reproduce the uniform operators bitwise. The released
stiffness is the joint weighted elimination of the free sides' traces and,
for a free-slip solid, of its parameters
(\testsym{test/indicator/geometry.jl}{Per-side free slip releases only the traces of its own sides}),
no corner trace is ever released
(\testsym{test/two_phase/variable_staggered_diffusion.jl}{LS-STAG per-side box traces follow their side and never a corner}),
a constant tangential velocity along a free-slip floor between hanging
vertices and viscosity jumps has no strain there
(\testsym{test/two_phase/variable_staggered_diffusion.jl}{Mixed per-side walls keep the rigid tangent of a free-slip side}),
and one along the sloping free-slip solid has none away from the box,
whatever the sides.

The manufactured streamfunctions vanish with their normal derivative on the
no-slip walls and with their second normal derivative on the free-slip
sides, whose tangential velocity is not zero. Between $n=64$ and $128$, the
orders of the shifted Helmholtz solve on uniform and right- or top-refined
meshes (LS-STAG) are
\begin{center}
\begin{tabular}{llcc}
\toprule
Sides & Solid & $L^2$ & Maximum \\
\midrule
free sides, no-slip floor, lid & none & $1.99$--$2.01$ & $1.94$--$1.98$ \\
free sides, no-slip lid & $y=c+sx$ meets free sides & $1.99$--$2.03$ & $1.74$--$1.96$ \\
free sides, no-slip lid & $y=sx$ through $(0,0)$ & $2.00$ & $1.96$--$1.97$ \\
no-slip sides, free lid & $y=c+sx$ & $2.00$--$2.06$ & $1.92$--$1.98$ \\
\bottomrule
\end{tabular}
\end{center}
The $1.74$ is a maximum in the interior, on the top-refined mesh where the
transition crosses the solid; it is $1.84$ between $n=128$ and $256$, while
the maximum within five cells of the box converges at $1.99$. Four of these
configurations, the corner one included, are pinned in
\testsym{test/two_phase/variable_staggered_diffusion.jl}{Mixed per-side walls keep second-order manufactured solutions}.

\paragraph{Cost of a thin Cartesian free-slip gap.}
If a wall runs parallel, or nearly parallel, to a box side only one to
two cells away, successive Green remainders may share traces along the
entire side. Eliminating them jointly is then nonlocal: on a patch with
$r$ stress rows and $f$ faces the stiffness contains a dense $f\times f$
block. This is a consequence of the exact minimum, not a cutoff in the
closure. With $S=W^{1/2}E_{\mathrm{patch}}$ after projection,
\srcsym{src/two_phase/diffusion/variable_diffusion_box_traces.jl}{_ls_add_box_stiffness!}
forms $S^\top S$ in a preallocated block and adds one triangle to both
orientations of the sparse stiffness. Ordinary stress rows retain their
separate sparse refresh map. For a strip with $r,f$ and trace rank all
proportional to the side length in cells, the stored patch data now grow
quadratically, not cubically. Dense arithmetic still costs cubic work
per refresh. The strip regression checks storage, nesting, and agreement
with fresh assembly after changing the viscosity:
\testsym{test/two_phase/variable_staggered_diffusion.jl}{LS-STAG thin-gap FreeSlip box traces are compact and refresh exactly}.

For a straight $1.5h$ fluid strip on a uniform $n\times n$ mesh (blocks of
$4\times4$, viscosity pair $2.5,0.4$), the table gives the complete retained
patch storage, warm construction time and cumulative construction allocation,
median of five viscosity refreshes, and stored stiffness entries. Julia
1.12.7, \texttt{Float64}, and one Julia and BLAS thread were used.
\begin{center}
\begin{tabular}{rlrrrrr}
\toprule
$n$ & Wall & Patch (MiB) & Build (ms) & Build (MiB) & Refresh (ms) & $\mathrm{nnz}(K)$ \\
\midrule
64  & No slip   & 0     & 4.04   & 9.97   & 0.081 & 9\,936 \\
64  & Free slip & 1.44  & 6.83   & 15.74  & 0.914 & 43\,974 \\
128 & No slip   & 0     & 15.48  & 40.29  & 0.315 & 36\,304 \\
128 & Free slip & 5.75  & 29.11  & 62.58  & 5.645 & 178\,054 \\
256 & No slip   & 0     & 107.06 & 154.33 & 1.317 & 138\,192 \\
256 & Free slip & 23.00 & 203.97 & 236.72 & 41.726 & 716\,550 \\
\bottomrule
\end{tabular}
\end{center}
The measured refresh allocated zero bytes in each configuration. At $n=256$ the
free-slip patch has 513 stress rows, 766 faces and trace rank 257;
enumerating the row/face/face products previously cost over 6\,GiB for
the tuples alone. The dense block is substantially cheaper, but it still
costs more to store and refresh than the no-slip stencil.

This PLIC extension to variable stress is implemented in
\srcsym{src/two_phase/diffusion/variable_diffusion_ls_stag.jl}{_ls_stress_rows!}.
One-velocity flow and both two-velocity projections support it with either
viscosity mixing rule. Adaptive topology changes rebuild the rows; coarse
levels with restricted geometry use the same closure.
The two-layer manufactured channel of
\testsym{test/operators/embedded_variable_wall_mms.jl}{two-phase embedded-wall manufactured Helmholtz (#106)}
has an aligned no-slip wall and continuous tangential traction across its
$2:1$ viscosity jump. From $n=16$ to $32$ to $64$, the volume-weighted
$L^2$ orders for LS-STAG are $1.88,1.97$ and $1.97,2.00$ for the two
components on the uniform mesh, and $1.95,1.99$ and $2.03,2.03$ with
the right half refined. This is not an order result for an oblique liquid
contact line.

\begin{sloppypar}
For the legacy wall-distance closure, let $q_{c,v}$ be the quarter of cell $c$ at its corner $v$, and let
$\mathcal W_r$ contain the quarters $(c,v)$ for which $c$ has a Cartesian wall
segment with $v$ on its solid side and the contact of row $r$ for the
wall-tangential component crosses the line through $v$ parallel to the wall:
that shear across the solid is supplied by the tangential wall rows (on a
uniform mesh these are the quarters at the wall corners of the cells). The
stress weights are the measures of the remaining rows times the viscosity
there,
\begin{equation}\label{eq:viscosity:weights}
  \begin{gathered}
  \stressw=\diag{2\meas{\opencell{c}}\visc_c,\ \ \meas{\shearregion{r}}^{\open}\visc_r,\ \ \wallcond{w}\visc_{c(w)},\ \ 1_{\mathrm{prescribed}}},\\
  \meas{\shearregion{r}}^{\open}=\sum_{(c,v)\notin\mathcal W_r}\openfrac{c}\,\meas{\shearregion{r}\cap q_{c,v}},\qquad
  \frac{1}{\visc_r}=\frac{1}{\meas{\shearregion{r}}^{\open}}\sum_{(c,v)\notin\mathcal W_r,\ \openfrac{c}>0}\openfrac{c}\,\meas{\shearregion{r}\cap q_{c,v}}\frac{1}{\visc_c},
  \end{gathered}
\end{equation}
so that $\visc_r$ is the open-area-weighted harmonic mean of the viscosities
of the cell parts the row covers (on a uniform mesh, the four quarters at a
vertex). This is the series conductance across those cells: the centred difference
\cref{eq:viscosity:shear} differences two velocities through cells that lie in
series, so their traction $\visc\partial_t\velo_k$ is transmitted exactly only
if the coefficient is the harmonic mean. For a layer aligned with the shear and
$\visc=(2,1)$ the row coefficient is $4/3$; the geometric mean
$\sqrt{2}$ overestimates the flux by $2\sqrt{\visc_1\visc_2}/(\visc_1+\visc_2)$,
$6.07\,\%$ there, and the 1-D shear across such a layer converges with order
$0.8$--$1.0$ instead of $2$ (\issue{113}). With
\apitype{GeometricViscosityAveraging} the operator uses the thesis'
$\prod_c\visc_c^{w_c}$, the geometric mean at the open-volume fraction, and
reproduces that behaviour; both rules keep $\visc_r>0$, so the stiffness
remains symmetric positive semidefinite and the energy balance of
\cref{sec:viscosity:time} holds for either. Cartesian embedded-wall rows use the owned segment
length divided by the wall-normal distance as in the componentwise operator;
a non-Cartesian PLIC wall retains the earlier midpoint-distance conductance,
half the open fraction of the face times the segment length divided by the
distance from the face centre to the segment midpoint, bounded below by a
quarter of the smallest cell width. That closure, and the cut closure
of~\cref{eq:viscosity:cut}, are not continuous in the tilted-wall limit
(\issue{110}). Wall rows use the viscosity
of the cell of their solid segment; covered unknowns outside the constrained
wall neighbourhood retain their algebraic penalty and the viscosity of
$\cminus{f}$. The integrated stiffness and the operator returned by
\apifun{diffusion!} are
\begin{equation}\label{eq:viscosity:operator}
  \strainmat^{\transpose}\stressw\strainmat,
  \qquad
  -\dualvolmat^{-1}\strainmat^{\transpose}\stressw\strainmat\,\velo,
\end{equation}
with $\dualvolmat=\diag{\dualvol{f}}$, or $\diag{\opendualvol{f}}$ (one where
it vanishes) with a solid; \apifun{diffusion!} returns zero at prescribed
faces. The stiffness is symmetric positive
semidefinite, and its quadratic form is the discrete dissipation
\begin{equation}\label{eq:viscosity:quadratic}
  \velo^{\transpose}\strainmat^{\transpose}\stressw\strainmat\velo
  =\sum_{c,k}2\visc_c\meas{\opencell{c}}(\partial_k\velo_k)_c^2
  +\sum_{r\in\shearrows}\visc_r\meas{\shearregion{r}}^{\open}\roundpar{(\vderiv{2}\velo_1)_{g(r)}+(\vderiv{1}\velo_2)_{g'(r)}}^2
  +\sum_{w}\wallcond{w}\visc_{c(w)}\velo_{f(w)}^2
  +\sum_{f\in\mathrm{prescribed}}\velo_f^2 ,
\end{equation}
the counterpart of $\integral{}{2\visc\strainrate:\strainrate}{V}$
in~\cref{eq:viscosity:dissipation}, which it approximates to second order
for a constant viscosity, also on composite meshes (derivation audit,
\code{proof_review/eq-viscosity-quadratic/}); for constant viscosity and
$\divergence\Velo=0$ the operator approximates $\visc\Delta\Velo$ to second
order on a uniform mesh (\code{proof_review/eq-viscosity-operator/}). At a
level transition the truncation error is $O(1)$ but in the flux
form~\cref{eq:viscosity:shear-flux} with first-order fluxes, and the solution
converges at second order (\issue{112}). The
kernel of the stiffness consists of the constant components on a periodic
mesh and of the constant tangential component in a channel that is periodic
along free-slip walls; a no-slip wall or a solid makes it definite. The
walls are part of the rows: no-slip
domain walls through the one-sided shear differences, free slip through the
vanishing ones, the embedded solid through the wall rows (no-slip only).
Only stationary walls are accepted, \apitype{NoSlip}\code{()} and
\apitype{FreeSlip}\code{()}, uniform or per side with
\apitype{VelocityBoundaryConditions} (\cref{sec:viscosity:per-side}); moving
walls and open sides are rejected. When the interface moves only $\stressw$ changes: the code
refreshes $\visc_c$ from the transported fractions before every viscous step
and keeps $\strainmat$ until the next regrid, which rebuilds the contacts and
rows from the lattice topology.
\end{sloppypar}

% Figure: rows of the strain matrix E near a level transition and a no-slip
% wall (Chapter 9). Cells, faces, the rows of E that contain one fine
% segment and the shear row at a hanging vertex are read from the
% code's VariableStaggeredDiffusionOperator: generated/viscosity.tex
% (viscosity.jl, viscosity_data). Labels hand-written.
\begin{figure}
  \centering
  \def\fxwnone{%
    \def\fxwcell##1##2##3##4##5{}\def\fxwface##1##2##3##4##5{}%
    \def\fxwrow##1##2##3##4{}\def\fxwstencil##1##2##3##4{}\def\fxwhanging##1##2{}}%
  % Background of both panels: level shading, cells, wall, faces (red:
  % normal e_1, blue: e_2; the target segment of panel (a) thick).
  \def\fxwbackground#1{%
    \begin{scope}
      \clip (3.0,0.0) rectangle (5.0,1.25);
      \fxwnone
      \def\fxwcell##1##2##3##4##5{\fill[tr level ##5] (##1,##2) rectangle (##3,##4);}%
      % Generated by docs/techref/figures/generate.jl; do not edit.
\fxwcell{3}{0}{3.5}{0.5}{0}
\fxwcell{3.5}{0}{4}{0.5}{0}
\fxwcell{3}{0.5}{3.5}{1}{0}
\fxwcell{3.5}{0.5}{4}{1}{0}
\fxwcell{3}{1}{3.5}{1.5}{0}
\fxwcell{3.5}{1}{4}{1.5}{0}
\fxwcell{4}{0}{4.25}{0.25}{1}
\fxwcell{4.25}{0}{4.5}{0.25}{1}
\fxwcell{4.5}{0}{4.75}{0.25}{1}
\fxwcell{4.75}{0}{5}{0.25}{1}
\fxwcell{4}{0.25}{4.25}{0.5}{1}
\fxwcell{4.25}{0.25}{4.5}{0.5}{1}
\fxwcell{4.5}{0.25}{4.75}{0.5}{1}
\fxwcell{4.75}{0.25}{5}{0.5}{1}
\fxwcell{4}{0.5}{4.25}{0.75}{1}
\fxwcell{4.25}{0.5}{4.5}{0.75}{1}
\fxwcell{4.5}{0.5}{4.75}{0.75}{1}
\fxwcell{4.75}{0.5}{5}{0.75}{1}
\fxwcell{4}{0.75}{4.25}{1}{1}
\fxwcell{4.25}{0.75}{4.5}{1}{1}
\fxwcell{4.5}{0.75}{4.75}{1}{1}
\fxwcell{4.75}{0.75}{5}{1}{1}
\fxwcell{4}{1}{4.25}{1.25}{1}
\fxwcell{4.25}{1}{4.5}{1.25}{1}
\fxwcell{4.5}{1}{4.75}{1.25}{1}
\fxwcell{4.75}{1}{5}{1.25}{1}
\fxwface{1}{3}{0.25}{0.5}{0}
\fxwface{1}{3}{0.75}{0.5}{0}
\fxwface{1}{3}{1.25}{0.5}{0}
\fxwface{1}{3.5}{0.25}{0.5}{0}
\fxwface{1}{3.5}{0.75}{0.5}{0}
\fxwface{1}{3.5}{1.25}{0.5}{0}
\fxwface{1}{4}{0.125}{0.25}{1}
\fxwface{1}{4}{0.375}{0.25}{0}
\fxwface{1}{4}{0.625}{0.25}{0}
\fxwface{1}{4}{0.875}{0.25}{0}
\fxwface{1}{4}{1.125}{0.25}{0}
\fxwface{1}{4.25}{0.125}{0.25}{0}
\fxwface{1}{4.25}{0.375}{0.25}{0}
\fxwface{1}{4.25}{0.625}{0.25}{0}
\fxwface{1}{4.25}{0.875}{0.25}{0}
\fxwface{1}{4.25}{1.125}{0.25}{0}
\fxwface{1}{4.5}{0.125}{0.25}{0}
\fxwface{1}{4.5}{0.375}{0.25}{0}
\fxwface{1}{4.5}{0.625}{0.25}{0}
\fxwface{1}{4.5}{0.875}{0.25}{0}
\fxwface{1}{4.5}{1.125}{0.25}{0}
\fxwface{1}{4.75}{0.125}{0.25}{0}
\fxwface{1}{4.75}{0.375}{0.25}{0}
\fxwface{1}{4.75}{0.625}{0.25}{0}
\fxwface{1}{4.75}{0.875}{0.25}{0}
\fxwface{1}{4.75}{1.125}{0.25}{0}
\fxwface{1}{5}{0.125}{0.25}{0}
\fxwface{1}{5}{0.375}{0.25}{0}
\fxwface{1}{5}{0.625}{0.25}{0}
\fxwface{1}{5}{0.875}{0.25}{0}
\fxwface{1}{5}{1.125}{0.25}{0}
\fxwface{2}{4.125}{0.25}{0.25}{0}
\fxwface{2}{4.375}{0.25}{0.25}{0}
\fxwface{2}{4.625}{0.25}{0.25}{0}
\fxwface{2}{4.875}{0.25}{0.25}{0}
\fxwface{2}{3.25}{0.5}{0.5}{0}
\fxwface{2}{3.75}{0.5}{0.5}{0}
\fxwface{2}{4.125}{0.5}{0.25}{0}
\fxwface{2}{4.375}{0.5}{0.25}{0}
\fxwface{2}{4.625}{0.5}{0.25}{0}
\fxwface{2}{4.875}{0.5}{0.25}{0}
\fxwface{2}{4.125}{0.75}{0.25}{0}
\fxwface{2}{4.375}{0.75}{0.25}{0}
\fxwface{2}{4.625}{0.75}{0.25}{0}
\fxwface{2}{4.875}{0.75}{0.25}{0}
\fxwface{2}{3.25}{1}{0.5}{0}
\fxwface{2}{3.75}{1}{0.5}{0}
\fxwface{2}{4.125}{1}{0.25}{0}
\fxwface{2}{4.375}{1}{0.25}{0}
\fxwface{2}{4.625}{1}{0.25}{0}
\fxwface{2}{4.875}{1}{0.25}{0}
\fxwface{2}{4.125}{1.25}{0.25}{0}
\fxwface{2}{4.375}{1.25}{0.25}{0}
\fxwface{2}{4.625}{1.25}{0.25}{0}
\fxwface{2}{4.875}{1.25}{0.25}{0}
\fxwrow{0}{3.75}{0.25}{1}
\fxwrow{0}{4.125}{0.125}{-4}
\fxwrow{1}{3.9375}{0.0625}{8}
\fxwrow{1}{3.9375}{0.3125}{-4}
\fxwrow{1}{3.9375}{0.1875}{-4}
\fxwhanging{4}{0.25}
\fxwstencil{1}{4}{0.125}{-4}
\fxwstencil{1}{4}{0.375}{4}
\fxwstencil{2}{4.125}{0.25}{2.67}
\fxwstencil{2}{3.75}{0.5}{-1.33}
%
      \fxwnone
      \def\fxwcell##1##2##3##4##5{\draw[tr cell] (##1,##2) rectangle (##3,##4);}%
      \def\fxwface##1##2##3##4##5{%
        \ifnum##5=1 \ifnum#1=1
          \draw[line width=2.4pt,trred] (##2,##3-##4/2+0.012) -- (##2,##3+##4/2-0.012);\fi\fi
        \ifnum##1=1 \draw[trred!80!black,line width=0.5pt,-{Stealth[length=2.4pt,width=2pt]}] (##2-0.035,##3) -- (##2+0.035,##3);
        \else \draw[trblue!80!black,line width=0.5pt,-{Stealth[length=2.4pt,width=2pt]}] (##2,##3-0.035) -- (##2,##3+0.035);\fi}%
      % Generated by docs/techref/figures/generate.jl; do not edit.
\fxwcell{3}{0}{3.5}{0.5}{0}
\fxwcell{3.5}{0}{4}{0.5}{0}
\fxwcell{3}{0.5}{3.5}{1}{0}
\fxwcell{3.5}{0.5}{4}{1}{0}
\fxwcell{3}{1}{3.5}{1.5}{0}
\fxwcell{3.5}{1}{4}{1.5}{0}
\fxwcell{4}{0}{4.25}{0.25}{1}
\fxwcell{4.25}{0}{4.5}{0.25}{1}
\fxwcell{4.5}{0}{4.75}{0.25}{1}
\fxwcell{4.75}{0}{5}{0.25}{1}
\fxwcell{4}{0.25}{4.25}{0.5}{1}
\fxwcell{4.25}{0.25}{4.5}{0.5}{1}
\fxwcell{4.5}{0.25}{4.75}{0.5}{1}
\fxwcell{4.75}{0.25}{5}{0.5}{1}
\fxwcell{4}{0.5}{4.25}{0.75}{1}
\fxwcell{4.25}{0.5}{4.5}{0.75}{1}
\fxwcell{4.5}{0.5}{4.75}{0.75}{1}
\fxwcell{4.75}{0.5}{5}{0.75}{1}
\fxwcell{4}{0.75}{4.25}{1}{1}
\fxwcell{4.25}{0.75}{4.5}{1}{1}
\fxwcell{4.5}{0.75}{4.75}{1}{1}
\fxwcell{4.75}{0.75}{5}{1}{1}
\fxwcell{4}{1}{4.25}{1.25}{1}
\fxwcell{4.25}{1}{4.5}{1.25}{1}
\fxwcell{4.5}{1}{4.75}{1.25}{1}
\fxwcell{4.75}{1}{5}{1.25}{1}
\fxwface{1}{3}{0.25}{0.5}{0}
\fxwface{1}{3}{0.75}{0.5}{0}
\fxwface{1}{3}{1.25}{0.5}{0}
\fxwface{1}{3.5}{0.25}{0.5}{0}
\fxwface{1}{3.5}{0.75}{0.5}{0}
\fxwface{1}{3.5}{1.25}{0.5}{0}
\fxwface{1}{4}{0.125}{0.25}{1}
\fxwface{1}{4}{0.375}{0.25}{0}
\fxwface{1}{4}{0.625}{0.25}{0}
\fxwface{1}{4}{0.875}{0.25}{0}
\fxwface{1}{4}{1.125}{0.25}{0}
\fxwface{1}{4.25}{0.125}{0.25}{0}
\fxwface{1}{4.25}{0.375}{0.25}{0}
\fxwface{1}{4.25}{0.625}{0.25}{0}
\fxwface{1}{4.25}{0.875}{0.25}{0}
\fxwface{1}{4.25}{1.125}{0.25}{0}
\fxwface{1}{4.5}{0.125}{0.25}{0}
\fxwface{1}{4.5}{0.375}{0.25}{0}
\fxwface{1}{4.5}{0.625}{0.25}{0}
\fxwface{1}{4.5}{0.875}{0.25}{0}
\fxwface{1}{4.5}{1.125}{0.25}{0}
\fxwface{1}{4.75}{0.125}{0.25}{0}
\fxwface{1}{4.75}{0.375}{0.25}{0}
\fxwface{1}{4.75}{0.625}{0.25}{0}
\fxwface{1}{4.75}{0.875}{0.25}{0}
\fxwface{1}{4.75}{1.125}{0.25}{0}
\fxwface{1}{5}{0.125}{0.25}{0}
\fxwface{1}{5}{0.375}{0.25}{0}
\fxwface{1}{5}{0.625}{0.25}{0}
\fxwface{1}{5}{0.875}{0.25}{0}
\fxwface{1}{5}{1.125}{0.25}{0}
\fxwface{2}{4.125}{0.25}{0.25}{0}
\fxwface{2}{4.375}{0.25}{0.25}{0}
\fxwface{2}{4.625}{0.25}{0.25}{0}
\fxwface{2}{4.875}{0.25}{0.25}{0}
\fxwface{2}{3.25}{0.5}{0.5}{0}
\fxwface{2}{3.75}{0.5}{0.5}{0}
\fxwface{2}{4.125}{0.5}{0.25}{0}
\fxwface{2}{4.375}{0.5}{0.25}{0}
\fxwface{2}{4.625}{0.5}{0.25}{0}
\fxwface{2}{4.875}{0.5}{0.25}{0}
\fxwface{2}{4.125}{0.75}{0.25}{0}
\fxwface{2}{4.375}{0.75}{0.25}{0}
\fxwface{2}{4.625}{0.75}{0.25}{0}
\fxwface{2}{4.875}{0.75}{0.25}{0}
\fxwface{2}{3.25}{1}{0.5}{0}
\fxwface{2}{3.75}{1}{0.5}{0}
\fxwface{2}{4.125}{1}{0.25}{0}
\fxwface{2}{4.375}{1}{0.25}{0}
\fxwface{2}{4.625}{1}{0.25}{0}
\fxwface{2}{4.875}{1}{0.25}{0}
\fxwface{2}{4.125}{1.25}{0.25}{0}
\fxwface{2}{4.375}{1.25}{0.25}{0}
\fxwface{2}{4.625}{1.25}{0.25}{0}
\fxwface{2}{4.875}{1.25}{0.25}{0}
\fxwrow{0}{3.75}{0.25}{1}
\fxwrow{0}{4.125}{0.125}{-4}
\fxwrow{1}{3.9375}{0.0625}{8}
\fxwrow{1}{3.9375}{0.3125}{-4}
\fxwrow{1}{3.9375}{0.1875}{-4}
\fxwhanging{4}{0.25}
\fxwstencil{1}{4}{0.125}{-4}
\fxwstencil{1}{4}{0.375}{4}
\fxwstencil{2}{4.125}{0.25}{2.67}
\fxwstencil{2}{3.75}{0.5}{-1.33}
%
    \end{scope}
    \fill[tr wall] (3.0,-0.1) rectangle (5.0,0.0);
    \draw[black,line width=0.8pt] (3.0,0.0) -- (5.0,0.0);
    \draw[black!70,dash pattern=on 2pt off 1.5pt,line width=0.5pt] (4,-0.12) -- (4,1.3);
    \fxwnone
    \def\fxwhanging##1##2{\node[tr vertex,minimum size=3.6pt] at (##1,##2) {};}%
    % Generated by docs/techref/figures/generate.jl; do not edit.
\fxwcell{3}{0}{3.5}{0.5}{0}
\fxwcell{3.5}{0}{4}{0.5}{0}
\fxwcell{3}{0.5}{3.5}{1}{0}
\fxwcell{3.5}{0.5}{4}{1}{0}
\fxwcell{3}{1}{3.5}{1.5}{0}
\fxwcell{3.5}{1}{4}{1.5}{0}
\fxwcell{4}{0}{4.25}{0.25}{1}
\fxwcell{4.25}{0}{4.5}{0.25}{1}
\fxwcell{4.5}{0}{4.75}{0.25}{1}
\fxwcell{4.75}{0}{5}{0.25}{1}
\fxwcell{4}{0.25}{4.25}{0.5}{1}
\fxwcell{4.25}{0.25}{4.5}{0.5}{1}
\fxwcell{4.5}{0.25}{4.75}{0.5}{1}
\fxwcell{4.75}{0.25}{5}{0.5}{1}
\fxwcell{4}{0.5}{4.25}{0.75}{1}
\fxwcell{4.25}{0.5}{4.5}{0.75}{1}
\fxwcell{4.5}{0.5}{4.75}{0.75}{1}
\fxwcell{4.75}{0.5}{5}{0.75}{1}
\fxwcell{4}{0.75}{4.25}{1}{1}
\fxwcell{4.25}{0.75}{4.5}{1}{1}
\fxwcell{4.5}{0.75}{4.75}{1}{1}
\fxwcell{4.75}{0.75}{5}{1}{1}
\fxwcell{4}{1}{4.25}{1.25}{1}
\fxwcell{4.25}{1}{4.5}{1.25}{1}
\fxwcell{4.5}{1}{4.75}{1.25}{1}
\fxwcell{4.75}{1}{5}{1.25}{1}
\fxwface{1}{3}{0.25}{0.5}{0}
\fxwface{1}{3}{0.75}{0.5}{0}
\fxwface{1}{3}{1.25}{0.5}{0}
\fxwface{1}{3.5}{0.25}{0.5}{0}
\fxwface{1}{3.5}{0.75}{0.5}{0}
\fxwface{1}{3.5}{1.25}{0.5}{0}
\fxwface{1}{4}{0.125}{0.25}{1}
\fxwface{1}{4}{0.375}{0.25}{0}
\fxwface{1}{4}{0.625}{0.25}{0}
\fxwface{1}{4}{0.875}{0.25}{0}
\fxwface{1}{4}{1.125}{0.25}{0}
\fxwface{1}{4.25}{0.125}{0.25}{0}
\fxwface{1}{4.25}{0.375}{0.25}{0}
\fxwface{1}{4.25}{0.625}{0.25}{0}
\fxwface{1}{4.25}{0.875}{0.25}{0}
\fxwface{1}{4.25}{1.125}{0.25}{0}
\fxwface{1}{4.5}{0.125}{0.25}{0}
\fxwface{1}{4.5}{0.375}{0.25}{0}
\fxwface{1}{4.5}{0.625}{0.25}{0}
\fxwface{1}{4.5}{0.875}{0.25}{0}
\fxwface{1}{4.5}{1.125}{0.25}{0}
\fxwface{1}{4.75}{0.125}{0.25}{0}
\fxwface{1}{4.75}{0.375}{0.25}{0}
\fxwface{1}{4.75}{0.625}{0.25}{0}
\fxwface{1}{4.75}{0.875}{0.25}{0}
\fxwface{1}{4.75}{1.125}{0.25}{0}
\fxwface{1}{5}{0.125}{0.25}{0}
\fxwface{1}{5}{0.375}{0.25}{0}
\fxwface{1}{5}{0.625}{0.25}{0}
\fxwface{1}{5}{0.875}{0.25}{0}
\fxwface{1}{5}{1.125}{0.25}{0}
\fxwface{2}{4.125}{0.25}{0.25}{0}
\fxwface{2}{4.375}{0.25}{0.25}{0}
\fxwface{2}{4.625}{0.25}{0.25}{0}
\fxwface{2}{4.875}{0.25}{0.25}{0}
\fxwface{2}{3.25}{0.5}{0.5}{0}
\fxwface{2}{3.75}{0.5}{0.5}{0}
\fxwface{2}{4.125}{0.5}{0.25}{0}
\fxwface{2}{4.375}{0.5}{0.25}{0}
\fxwface{2}{4.625}{0.5}{0.25}{0}
\fxwface{2}{4.875}{0.5}{0.25}{0}
\fxwface{2}{4.125}{0.75}{0.25}{0}
\fxwface{2}{4.375}{0.75}{0.25}{0}
\fxwface{2}{4.625}{0.75}{0.25}{0}
\fxwface{2}{4.875}{0.75}{0.25}{0}
\fxwface{2}{3.25}{1}{0.5}{0}
\fxwface{2}{3.75}{1}{0.5}{0}
\fxwface{2}{4.125}{1}{0.25}{0}
\fxwface{2}{4.375}{1}{0.25}{0}
\fxwface{2}{4.625}{1}{0.25}{0}
\fxwface{2}{4.875}{1}{0.25}{0}
\fxwface{2}{4.125}{1.25}{0.25}{0}
\fxwface{2}{4.375}{1.25}{0.25}{0}
\fxwface{2}{4.625}{1.25}{0.25}{0}
\fxwface{2}{4.875}{1.25}{0.25}{0}
\fxwrow{0}{3.75}{0.25}{1}
\fxwrow{0}{4.125}{0.125}{-4}
\fxwrow{1}{3.9375}{0.0625}{8}
\fxwrow{1}{3.9375}{0.3125}{-4}
\fxwrow{1}{3.9375}{0.1875}{-4}
\fxwhanging{4}{0.25}
\fxwstencil{1}{4}{0.125}{-4}
\fxwstencil{1}{4}{0.375}{4}
\fxwstencil{2}{4.125}{0.25}{2.67}
\fxwstencil{2}{3.75}{0.5}{-1.33}
%
  }%
  \begin{subfigure}[b]{0.495\textwidth}
    \centering
    \begin{tikzpicture}[x=3.05cm,y=3.05cm]
      \fxwbackground{1}
      \fxwnone
      \def\fxwrow#1#2#3#4{%
        \ifnum#1=0
          \node[rectangle,draw=black,fill=trgreen!55,inner sep=0pt,minimum size=4.4pt] at (#2,#3) {};
          \node[tr value,anchor=south west] at (#2+0.015,#3+0.015) {$#4$};
        \else
          \node[diamond,draw=black,fill=trpurple!55,inner sep=0pt,minimum size=5.6pt] at (#2,#3) {};
          \node[tr value,anchor=west] at (#2+0.03,#3+0.05) {$#4$};
        \fi}%
      % Generated by docs/techref/figures/generate.jl; do not edit.
\fxwcell{3}{0}{3.5}{0.5}{0}
\fxwcell{3.5}{0}{4}{0.5}{0}
\fxwcell{3}{0.5}{3.5}{1}{0}
\fxwcell{3.5}{0.5}{4}{1}{0}
\fxwcell{3}{1}{3.5}{1.5}{0}
\fxwcell{3.5}{1}{4}{1.5}{0}
\fxwcell{4}{0}{4.25}{0.25}{1}
\fxwcell{4.25}{0}{4.5}{0.25}{1}
\fxwcell{4.5}{0}{4.75}{0.25}{1}
\fxwcell{4.75}{0}{5}{0.25}{1}
\fxwcell{4}{0.25}{4.25}{0.5}{1}
\fxwcell{4.25}{0.25}{4.5}{0.5}{1}
\fxwcell{4.5}{0.25}{4.75}{0.5}{1}
\fxwcell{4.75}{0.25}{5}{0.5}{1}
\fxwcell{4}{0.5}{4.25}{0.75}{1}
\fxwcell{4.25}{0.5}{4.5}{0.75}{1}
\fxwcell{4.5}{0.5}{4.75}{0.75}{1}
\fxwcell{4.75}{0.5}{5}{0.75}{1}
\fxwcell{4}{0.75}{4.25}{1}{1}
\fxwcell{4.25}{0.75}{4.5}{1}{1}
\fxwcell{4.5}{0.75}{4.75}{1}{1}
\fxwcell{4.75}{0.75}{5}{1}{1}
\fxwcell{4}{1}{4.25}{1.25}{1}
\fxwcell{4.25}{1}{4.5}{1.25}{1}
\fxwcell{4.5}{1}{4.75}{1.25}{1}
\fxwcell{4.75}{1}{5}{1.25}{1}
\fxwface{1}{3}{0.25}{0.5}{0}
\fxwface{1}{3}{0.75}{0.5}{0}
\fxwface{1}{3}{1.25}{0.5}{0}
\fxwface{1}{3.5}{0.25}{0.5}{0}
\fxwface{1}{3.5}{0.75}{0.5}{0}
\fxwface{1}{3.5}{1.25}{0.5}{0}
\fxwface{1}{4}{0.125}{0.25}{1}
\fxwface{1}{4}{0.375}{0.25}{0}
\fxwface{1}{4}{0.625}{0.25}{0}
\fxwface{1}{4}{0.875}{0.25}{0}
\fxwface{1}{4}{1.125}{0.25}{0}
\fxwface{1}{4.25}{0.125}{0.25}{0}
\fxwface{1}{4.25}{0.375}{0.25}{0}
\fxwface{1}{4.25}{0.625}{0.25}{0}
\fxwface{1}{4.25}{0.875}{0.25}{0}
\fxwface{1}{4.25}{1.125}{0.25}{0}
\fxwface{1}{4.5}{0.125}{0.25}{0}
\fxwface{1}{4.5}{0.375}{0.25}{0}
\fxwface{1}{4.5}{0.625}{0.25}{0}
\fxwface{1}{4.5}{0.875}{0.25}{0}
\fxwface{1}{4.5}{1.125}{0.25}{0}
\fxwface{1}{4.75}{0.125}{0.25}{0}
\fxwface{1}{4.75}{0.375}{0.25}{0}
\fxwface{1}{4.75}{0.625}{0.25}{0}
\fxwface{1}{4.75}{0.875}{0.25}{0}
\fxwface{1}{4.75}{1.125}{0.25}{0}
\fxwface{1}{5}{0.125}{0.25}{0}
\fxwface{1}{5}{0.375}{0.25}{0}
\fxwface{1}{5}{0.625}{0.25}{0}
\fxwface{1}{5}{0.875}{0.25}{0}
\fxwface{1}{5}{1.125}{0.25}{0}
\fxwface{2}{4.125}{0.25}{0.25}{0}
\fxwface{2}{4.375}{0.25}{0.25}{0}
\fxwface{2}{4.625}{0.25}{0.25}{0}
\fxwface{2}{4.875}{0.25}{0.25}{0}
\fxwface{2}{3.25}{0.5}{0.5}{0}
\fxwface{2}{3.75}{0.5}{0.5}{0}
\fxwface{2}{4.125}{0.5}{0.25}{0}
\fxwface{2}{4.375}{0.5}{0.25}{0}
\fxwface{2}{4.625}{0.5}{0.25}{0}
\fxwface{2}{4.875}{0.5}{0.25}{0}
\fxwface{2}{4.125}{0.75}{0.25}{0}
\fxwface{2}{4.375}{0.75}{0.25}{0}
\fxwface{2}{4.625}{0.75}{0.25}{0}
\fxwface{2}{4.875}{0.75}{0.25}{0}
\fxwface{2}{3.25}{1}{0.5}{0}
\fxwface{2}{3.75}{1}{0.5}{0}
\fxwface{2}{4.125}{1}{0.25}{0}
\fxwface{2}{4.375}{1}{0.25}{0}
\fxwface{2}{4.625}{1}{0.25}{0}
\fxwface{2}{4.875}{1}{0.25}{0}
\fxwface{2}{4.125}{1.25}{0.25}{0}
\fxwface{2}{4.375}{1.25}{0.25}{0}
\fxwface{2}{4.625}{1.25}{0.25}{0}
\fxwface{2}{4.875}{1.25}{0.25}{0}
\fxwrow{0}{3.75}{0.25}{1}
\fxwrow{0}{4.125}{0.125}{-4}
\fxwrow{1}{3.9375}{0.0625}{8}
\fxwrow{1}{3.9375}{0.3125}{-4}
\fxwrow{1}{3.9375}{0.1875}{-4}
\fxwhanging{4}{0.25}
\fxwstencil{1}{4}{0.125}{-4}
\fxwstencil{1}{4}{0.375}{4}
\fxwstencil{2}{4.125}{0.25}{2.67}
\fxwstencil{2}{3.75}{0.5}{-1.33}

      \node[tr small,anchor=north east,fill=white,fill opacity=0.85,text opacity=1] at (3.98,1.23) {$c$};
      \node[tr small,anchor=north west,fill=white,fill opacity=0.85,text opacity=1] at (4.27,0.23) {$c'$};
    \end{tikzpicture}
    \caption{Rows of $\strainmat$ that contain $\velo_f$}
    \label{fig:viscosity:stencil:rows}
  \end{subfigure}\hfill
  \begin{subfigure}[b]{0.495\textwidth}
    \centering
    \begin{tikzpicture}[x=3.05cm,y=3.05cm]
      \fxwbackground{2}
      \fxwnone
      \def\fxwstencil#1#2#3#4{%
        \ifnum#1=1 \node[circle,fill=trred,draw=black,line width=0.3pt,inner sep=0pt,minimum size=3.6pt] at (#2,#3) {};
        \else \node[circle,fill=trblue,draw=black,line width=0.3pt,inner sep=0pt,minimum size=3.6pt] at (#2,#3) {};\fi
        \node[tr value,anchor=south] at (#2,#3+0.03) {$#4$};}%
      % Generated by docs/techref/figures/generate.jl; do not edit.
\fxwcell{3}{0}{3.5}{0.5}{0}
\fxwcell{3.5}{0}{4}{0.5}{0}
\fxwcell{3}{0.5}{3.5}{1}{0}
\fxwcell{3.5}{0.5}{4}{1}{0}
\fxwcell{3}{1}{3.5}{1.5}{0}
\fxwcell{3.5}{1}{4}{1.5}{0}
\fxwcell{4}{0}{4.25}{0.25}{1}
\fxwcell{4.25}{0}{4.5}{0.25}{1}
\fxwcell{4.5}{0}{4.75}{0.25}{1}
\fxwcell{4.75}{0}{5}{0.25}{1}
\fxwcell{4}{0.25}{4.25}{0.5}{1}
\fxwcell{4.25}{0.25}{4.5}{0.5}{1}
\fxwcell{4.5}{0.25}{4.75}{0.5}{1}
\fxwcell{4.75}{0.25}{5}{0.5}{1}
\fxwcell{4}{0.5}{4.25}{0.75}{1}
\fxwcell{4.25}{0.5}{4.5}{0.75}{1}
\fxwcell{4.5}{0.5}{4.75}{0.75}{1}
\fxwcell{4.75}{0.5}{5}{0.75}{1}
\fxwcell{4}{0.75}{4.25}{1}{1}
\fxwcell{4.25}{0.75}{4.5}{1}{1}
\fxwcell{4.5}{0.75}{4.75}{1}{1}
\fxwcell{4.75}{0.75}{5}{1}{1}
\fxwcell{4}{1}{4.25}{1.25}{1}
\fxwcell{4.25}{1}{4.5}{1.25}{1}
\fxwcell{4.5}{1}{4.75}{1.25}{1}
\fxwcell{4.75}{1}{5}{1.25}{1}
\fxwface{1}{3}{0.25}{0.5}{0}
\fxwface{1}{3}{0.75}{0.5}{0}
\fxwface{1}{3}{1.25}{0.5}{0}
\fxwface{1}{3.5}{0.25}{0.5}{0}
\fxwface{1}{3.5}{0.75}{0.5}{0}
\fxwface{1}{3.5}{1.25}{0.5}{0}
\fxwface{1}{4}{0.125}{0.25}{1}
\fxwface{1}{4}{0.375}{0.25}{0}
\fxwface{1}{4}{0.625}{0.25}{0}
\fxwface{1}{4}{0.875}{0.25}{0}
\fxwface{1}{4}{1.125}{0.25}{0}
\fxwface{1}{4.25}{0.125}{0.25}{0}
\fxwface{1}{4.25}{0.375}{0.25}{0}
\fxwface{1}{4.25}{0.625}{0.25}{0}
\fxwface{1}{4.25}{0.875}{0.25}{0}
\fxwface{1}{4.25}{1.125}{0.25}{0}
\fxwface{1}{4.5}{0.125}{0.25}{0}
\fxwface{1}{4.5}{0.375}{0.25}{0}
\fxwface{1}{4.5}{0.625}{0.25}{0}
\fxwface{1}{4.5}{0.875}{0.25}{0}
\fxwface{1}{4.5}{1.125}{0.25}{0}
\fxwface{1}{4.75}{0.125}{0.25}{0}
\fxwface{1}{4.75}{0.375}{0.25}{0}
\fxwface{1}{4.75}{0.625}{0.25}{0}
\fxwface{1}{4.75}{0.875}{0.25}{0}
\fxwface{1}{4.75}{1.125}{0.25}{0}
\fxwface{1}{5}{0.125}{0.25}{0}
\fxwface{1}{5}{0.375}{0.25}{0}
\fxwface{1}{5}{0.625}{0.25}{0}
\fxwface{1}{5}{0.875}{0.25}{0}
\fxwface{1}{5}{1.125}{0.25}{0}
\fxwface{2}{4.125}{0.25}{0.25}{0}
\fxwface{2}{4.375}{0.25}{0.25}{0}
\fxwface{2}{4.625}{0.25}{0.25}{0}
\fxwface{2}{4.875}{0.25}{0.25}{0}
\fxwface{2}{3.25}{0.5}{0.5}{0}
\fxwface{2}{3.75}{0.5}{0.5}{0}
\fxwface{2}{4.125}{0.5}{0.25}{0}
\fxwface{2}{4.375}{0.5}{0.25}{0}
\fxwface{2}{4.625}{0.5}{0.25}{0}
\fxwface{2}{4.875}{0.5}{0.25}{0}
\fxwface{2}{4.125}{0.75}{0.25}{0}
\fxwface{2}{4.375}{0.75}{0.25}{0}
\fxwface{2}{4.625}{0.75}{0.25}{0}
\fxwface{2}{4.875}{0.75}{0.25}{0}
\fxwface{2}{3.25}{1}{0.5}{0}
\fxwface{2}{3.75}{1}{0.5}{0}
\fxwface{2}{4.125}{1}{0.25}{0}
\fxwface{2}{4.375}{1}{0.25}{0}
\fxwface{2}{4.625}{1}{0.25}{0}
\fxwface{2}{4.875}{1}{0.25}{0}
\fxwface{2}{4.125}{1.25}{0.25}{0}
\fxwface{2}{4.375}{1.25}{0.25}{0}
\fxwface{2}{4.625}{1.25}{0.25}{0}
\fxwface{2}{4.875}{1.25}{0.25}{0}
\fxwrow{0}{3.75}{0.25}{1}
\fxwrow{0}{4.125}{0.125}{-4}
\fxwrow{1}{3.9375}{0.0625}{8}
\fxwrow{1}{3.9375}{0.3125}{-4}
\fxwrow{1}{3.9375}{0.1875}{-4}
\fxwhanging{4}{0.25}
\fxwstencil{1}{4}{0.125}{-4}
\fxwstencil{1}{4}{0.375}{4}
\fxwstencil{2}{4.125}{0.25}{2.67}
\fxwstencil{2}{3.75}{0.5}{-1.33}

    \end{tikzpicture}
    \caption{Shear rows at the hanging vertex}
    \label{fig:viscosity:stencil:shear}
  \end{subfigure}
  \caption{The strain-rate rows~\cref{eq:viscosity:rows} next to a level
    transition (dashed) and a no-slip wall (hatched), read from the code's
    operator on the walled composite fixture (coarse cells of width $\half$
    on the left). (a)~The \chkViscFigNormal\ normal rows (green squares, at
    the centres of the coarse cell $c$ and the fine cell $c'$) and the
    \chkViscFigShear\ shear rows (purple diamonds, at the centres of their
    regions~\cref{eq:viscosity:shear-partition}) whose entry for the fine
    segment $f$ (thick red) is not zero, with that entry:
    $\pm\aperture{f}\meas{f}/\meas{\opencell{c}}$ in the normal rows, the
    two-point difference to the wall and across the hanging vertex in the
    shear rows. (b)~The two shear rows that meet at the hanging vertex (open
    circle) have the same \chkViscFigStencil\ unknowns: $\velo_1$ (red) by the
    two-point difference along the fine segments, $\velo_2$ (blue) by the
    anchored derivative of~\cref{eq:viscosity:shear} of the fine segment
    through the vertex, between the coarse unknown above and the zero wall
    value below, since the coarse cell has no $\velo_2$ unknown on the line
    through the vertex.}
  \label{fig:viscosity:stencil}
\end{figure}


\emph{Checks} (\code{stress_rows_check}, \code{affine_strain_check}).
The 1-D shear $\visc(y)\velo_1'(y)=-\cos(2\pi y)$ across an interface aligned
with $y=1/2$, whose tangential traction is continuous and whose velocity
derivative jumps by the viscosity ratio, converges at order $2.03$, $2.01$,
$2.00$ and $2.00$ ($\visc=(2,1)$, $n=16\to128$; the equal-viscosity solve gives
the same orders), while \apitype{GeometricViscosityAveraging} reproduces the
reported first-order rates below $1.2$
(\testsym{test/two_phase/variable_staggered_diffusion.jl}{Aligned viscosity jump transmits the tangential traction (#113)},
\testsym{test/two_phase/variable_staggered_diffusion.jl}{Geometric mixture keeps the first-order aligned shear (#113)}).
Both rules keep the stiffness symmetric and the implicit midpoint step
dissipative on a composite mesh with an aligned jump
(\testsym{test/two_phase/variable_staggered_diffusion.jl}{Harmonic vertex viscosity stays dissipative (#113)}). An
independent assembly of the normal rows~\cref{eq:viscosity:rows} equals the
code's to $\chkViscNormalDiff$ on the composite, cut and periodic composite
drop fixtures ($\visc=(0.1,0.002)$); the independent
weights~\cref{eq:viscosity:weights} of all normal and \chkViscShearRows\
shear rows (composite fixture) agree to $\chkViscWeightDiff$ (relative); the
cut fixture has \chkViscWallRows\ wall rows. For the six affine velocities
and both constants, the \chkViscAffineRows\ rows more than four coarse cells
from the domain edges return the exact strain rates to $\chkViscAffine$,
on the walled and on the periodic composite mesh.
On fixed composite trees (\issue{112}) the affine velocities give exact shear
rates and vanishing shear forces on every staggered volume at least two coarse
cells from the walls and periodic seams, with no-slip and free-slip walls,
the rows partition the periodic domain, the stiffness is symmetric positive
semidefinite and its kernel contains the unconstrained constant components
(\testsym{test/two_phase/variable_staggered_diffusion.jl}{Variable stress shear rows are flux-consistent at level transitions (#112)}).
The shifted solve $(\dualvolmat+0.01\,\strainmat^{\transpose}\stressw\strainmat)\velo=\dots$
for a manufactured no-slip velocity with $\visc=1+\tfrac14\cos(2\pi x)$ on the
walled tree with a doubly refined corner block converges with the orders
$2.01$, $2.01$ in $L^2$ and $2.00$, $1.95$ in the maximum norm ($n=8\to32$
cells per block side; before \issue{112}: $0.99$, $1.01$ in $L^2$ and $0.62$,
$0.80$ in the maximum norm), both for the fixed and the adapted tree
(\testsym{test/two_phase/variable_staggered_diffusion.jl}{Variable stress shifted solve on fixed and adapted trees}).
In a wider study (constant and variable viscosity, periodic and walled,
corner and strip refinement, $n=8\to64$) the $L^2$ orders are $1.96$--$2.03$
and the maximum-norm orders $1.53$--$2.00$, rising to $1.84$--$2.00$ at the
finest pair; the largest truncation error within three fine cells of a
transition stays between $3$ and $19$, where it grew as $h^{-1}$ before
(from $1.6\cdot10^2$--$7.2\cdot10^2$ at $n=8$ to $6.5\cdot10^2$--$3.0\cdot10^3$
at $n=32$).

\begin{deviation}{Two-phase stress rows}\label{dev:viscosity:stress}
The thesis locates the off-diagonal strain at the edges of the cells and
needs no interpolation on a rectilinear mesh, with the viscosity from the
volume fractions interpolated to the edges~\thesis{eq.~(7.52), \S7.5.1}. The
code uses the conditional fractions and places the shear rows on the
intersections of the contact regions of the two velocity components
(\cref{eq:viscosity:shear-partition}); on a conforming mesh these are the
vertex-centred dual cells with the open quarter-cell volumes, at level
transitions they keep the contact fluxes of both components consistent,
with the anchored fine--coarse derivatives of~\cref{eq:viscosity:shear}
(\issue{112}). At stationary solids, LS-STAG combines shared vertex shear
with the Green completion of \cref{eq:viscosity:lsstress}; the legacy closure
adds wall-distance rows instead. The cell and row mixtures are harmonic (series) by default,
\cref{eq:viscosity:viscosity,eq:viscosity:weights}; the thesis'
volume-fraction-weighted geometric mean remains available as
\apitype{GeometricViscosityAveraging} through the \code{viscosity_averaging}
option of \apitype{NavierStokesOptions}, because it is a diffuse model whose
row coefficient does not transmit the tangential traction across an aligned
interface at second order (\issue{113}).
\end{deviation}

\begin{implnote}
\apitype{VariableStaggeredDiffusionOperator}, built by
\apitype{StaggeredDiffusionOperator}\code{(interface; viscosity=(mu_l, mu_g), boundary=NoSlip())}
in \srcsym{src/two_phase/diffusion/variable_diffusion_operator.jl}{_variable_staggered_diffusion}~\cref{eq:viscosity:rows,eq:viscosity:weights};
\apifun{dynamic_viscosity} and \apitype{AbstractViscosityAveraging}
(\srcsym{src/operators/viscosity_averaging.jl}{HarmonicViscosityAveraging},
\srcsym{src/operators/viscosity_averaging.jl}{GeometricViscosityAveraging})
~\cref{eq:viscosity:viscosity}; stress-row mixtures
\srcsym{src/two_phase/diffusion/variable_diffusion_operator.jl}{_fill_stress_weights!}~\cref{eq:viscosity:weights};
shear partition \srcsym{src/two_phase/diffusion/variable_diffusion_geometry.jl}{_shear_contacts!}~\cref{eq:viscosity:shear-partition},
derivatives \srcsym{src/two_phase/diffusion/variable_diffusion_geometry.jl}{_shear_stencils!}~\cref{eq:viscosity:shear},
shear rows \srcsym{src/two_phase/diffusion/variable_diffusion_geometry.jl}{_shear_rows!}~\cref{eq:viscosity:rows,eq:viscosity:shear-flux};
embedded wall rows \srcsym{src/embedded/operators.jl}{_embedded_variable_wall_stresses};
stiffness \srcsym{src/two_phase/diffusion/variable_diffusion_geometry.jl}{_symmetric_stiffness!};
refresh \srcsym{src/two_phase/diffusion/variable_diffusion_operator.jl}{_update_dynamic_viscosity!};
action \srcsym{src/two_phase/diffusion/variable_diffusion_operator.jl}{_factorized_variable_stiffness_mul!}.
Tests: \testsym{test/two_phase/variable_staggered_diffusion.jl}{Factorized variable-stress action matches assembled operator},
\testsym{test/two_phase/variable_staggered_diffusion.jl}{Phase-viscosity averaging rules (#113)},
\testsym{test/two_phase/variable_staggered_diffusion.jl}{Variable coupled stress diffusion},
\testsym{test/two_phase/variable_staggered_diffusion.jl}{Variable stress manufactured bulk convergence},
\testsym{test/two_phase/variable_staggered_diffusion.jl}{variable diffusion is independent of base-root partitioning},
\testsym{test/two_phase/variable_staggered_diffusion.jl}{Variable stress shear rows are flux-consistent at level transitions (#112)},
\testsym{test/two_phase/variable_staggered_diffusion.jl}{Variable stress shifted solve on fixed and adapted trees},
\testsym{test/operators/embedded_variable_wall_mms.jl}{two-phase embedded-wall manufactured Helmholtz (#106)},
\testsym{test/operators/embedded_variable_wall_mms.jl}{two-phase aligned-wall constraints and reactions (#106)}.
\end{implnote}

\section{Fully discrete update: implicit midpoint solves}\label{sec:viscosity:time}
\emph{Single phase.} With $\velo^{\timelevel{n}}$ the old velocity with its
boundary values at $\timevar^{\timelevel{n}}$ and $\velo^{\stage{1}}$ the
transported velocity after the body acceleration and the boundary values at
$\timevar^{\timelevel{n}}+\half\dt$ (\cref{eq:forces:stages}), each component
solves
\begin{equation}\label{eq:viscosity:single-midpoint}
  \roundpar{\idmat-\tfrac12\dt\tfrac{\visc}{\dens}\lapstag{k}}\velo^{\stage{2}}
  =\velo^{\stage{1}}+\tfrac12\dt\tfrac{\visc}{\dens}\lapstag{k}\velo^{\timelevel{n}}
  +\dt\tfrac{\visc}{\dens}\dualvolmat^{-1}\sum_{w}\wallcond{w}\velo^{\physbnd}_w\big(\timevar^{\timelevel{n}}+\tfrac12\dt\big),
\end{equation}
a Crank--Nicolson step whose explicit half is evaluated at the old velocity
$\velo^{\timelevel{n}}$, not at the transported one, and whose wall data are
taken at the midpoint. The system is multiplied by $\dualvolmat$ and solved
by CG if $\dualvolmat\lapstag{k}$ is symmetric for every shift, including the
embedded conservative operator, otherwise by BiCGSTAB, both
preconditioned by default with the rediscretised tree multigrid of
\cref{ch:solvers}. The shift $\half\dt\visc/\dens$ of the system matrix, of every
multigrid level and of the bottom factorization is the current one: every
change of the step refreshes them
(\cref{sec:viscosity:refresh}).

\emph{Two phases, one velocity.} With the mass matrix
$\massmat=\diag{\stagdens_f\stagvol{f}}$ of the transported fractions
(\cref{eq:forces:density}; open volumes with a solid) and the viscosity
refreshed from the same fractions,
\begin{equation}\label{eq:viscosity:phase-midpoint}
  \roundpar{\massmat+\tfrac12\dt\,\strainmat^{\transpose}\stressw\strainmat}\velo^{\stage{2}}
  =\massmat\velo^{\stage{1}}-\tfrac12\dt\,\strainmat^{\transpose}\stressw\strainmat\,\velo^{\stage{1}},
\end{equation}
the implicit midpoint rule about the transported velocity, as in the
thesis' step T3.I with the stress~\thesis{eq.~(7.52)}; the matrix is
symmetric positive definite, solved by CG with the factorised action
$\massmat\velo+\half\dt\strainmat^{\transpose}(\stressw(\strainmat\velo))$
and the rediscretised multigrid of \cref{ch:solvers}. If both viscosities
are zero no operator is built and $\velo^{\stage{2}}=\velo^{\stage{1}}$;
one zero viscosity is rejected, as is a non-positive single-phase
viscosity.

\paragraph{Refresh on a step change.}\label{sec:viscosity:refresh}
The single-phase system of~\cref{eq:viscosity:single-midpoint} depends
on the step through the shift $\half\dt\visc/\dens$. A changed shift
refreshes the system and multigrid matrices and refactorises the bottom
system. The two-phase system of~\cref{eq:viscosity:phase-midpoint} also
depends on the evolving density and viscosity: its numerical hierarchy and
bottom factorization are refreshed every step even when $\dt$ is unchanged.
For the mimetic viscous solve, the band factor additionally depends on the
current metric, drag and step. Topology changes can also require rebuilding
these structures.

\emph{Stopping criterion.} With $\resvec$ the residual and $\rhsvec$ the
right-hand side of the integrated system and $\mathsf{m}_f$ the diagonal of
$\dualvolmat$ (single phase) or $\massmat$ (two phases), both solves stop
when $\max_f\abs{\resvec_f}/\mathsf{m}_f\le\max(\soltol\max_f\abs{\rhsvec_f}/\mathsf{m}_f,\ \text{round-off floor})$:
a velocity defect relative to the velocity scale of the right-hand side,
with the step tolerance $\soltol$ (default $\machineeps^{1/2}$) and at most
100 (one phase) or 500 (two phases) iterations; otherwise the step raises
an error.

\emph{Identities.} Let
$\resvec=\rhsvec-\roundpar{\massmat+\tfrac12\dt\,\strainmat^{\transpose}\stressw\strainmat}\velo^{\stage{2}}$
be the residual of~\cref{eq:viscosity:phase-midpoint}, with the right-hand
side $\rhsvec=\roundpar{\massmat-\tfrac12\dt\,\strainmat^{\transpose}\stressw\strainmat}\velo^{\stage{1}}$,
as in the stopping criterion. Then
$\massmat(\velo^{\stage{2}}-\velo^{\stage{1}})=-\tfrac12\dt\,\strainmat^{\transpose}\stressw\strainmat\,(\velo^{\stage{1}}+\velo^{\stage{2}})-\resvec$,
and multiplying by $\velo^{\stage{1}}+\velo^{\stage{2}}$ gives the energy
balance
\begin{equation}\label{eq:viscosity:energy}
  \tfrac12(\velo^{\stage{2}})^{\transpose}\massmat\velo^{\stage{2}}
  -\tfrac12(\velo^{\stage{1}})^{\transpose}\massmat\velo^{\stage{1}}
  =-\tfrac14\dt\,(\velo^{\stage{1}}+\velo^{\stage{2}})^{\transpose}\strainmat^{\transpose}\stressw\strainmat\,(\velo^{\stage{1}}+\velo^{\stage{2}})
  -\tfrac12(\velo^{\stage{1}}+\velo^{\stage{2}})^{\transpose}\resvec :
\end{equation}
the viscous step does not increase the mass-weighted kinetic energy
$\sum_f\half\stagdens_f\stagvol{f}\velo_f^2$, for every $\dt$, up to the
residual, the discrete counterpart of~\cref{eq:viscosity:dissipation}; the
residual term is at most
$\tfrac12\max_f(\abs{\resvec_f}/\mathsf{m}_f)\sum_f\mathsf{m}_f\abs{\velo^{\stage{1}}_f+\velo^{\stage{2}}_f}$,
with the maximum bounded by the stopping criterion.
The balance uses only that $\massmat$ is diagonal and positive and that the
stiffness is symmetric positive semidefinite, which
$\strainmat^{\transpose}\stressw\strainmat$ is by construction. It is an
equality: on an eigenvector of
$\massmat^{-1}\strainmat^{\transpose}\stressw\strainmat$ with eigenvalue
$\lambda$ an exact solve multiplies the energy by
$\roundpar{(1-\tfrac12\dt\lambda)/(1+\tfrac12\dt\lambda)}^2$, so the kernel
keeps its energy and stiff modes are hardly damped for large $\dt$ (the
midpoint rule is not L-stable). On a periodic mesh without solid a constant
component has no strain, so $\strainmat\faceone{k}=0$ for the indicator
$\faceone{k}$ of $\facesaxis{k}$ and, by symmetry,
\begin{equation}\label{eq:viscosity:momentum}
  \sum_{f\in\facesaxis{k}}\stagdens_f\stagvol{f}\velo^{\stage{2}}_f
  -\sum_{f\in\facesaxis{k}}\stagdens_f\stagvol{f}\velo^{\stage{1}}_f=-(\faceone{k})^{\transpose}\resvec :
\end{equation}
the momentum is conserved up to the residual. A solid exchanges momentum
through its wall rows; a wall exchanges the momentum of the component normal
to it through the normal rows, which omit the boundary faces (zero
wall-normal velocity), and a no-slip wall that of the tangential component
through the wall-adjacent shear differences, while a free-slip wall
conserves the tangential momentum on a mesh that is periodic along it. The
single-phase step~\cref{eq:viscosity:single-midpoint} conserves momentum on
any periodic composite mesh for any explicit argument, since
$(\faceone{k})^{\transpose}\dualvolmat\lapstag{k}=0$. Its backward-Euler
resolvent is contractive for every step size, and its midpoint update obeys
the same energy balance when $\velo^{\timelevel{n}}=\velo^{\stage{1}}$.
The actual split step evaluates the explicit half at
$\velo^{\timelevel{n}}$ and therefore has no such balance when these
velocities differ.

\emph{Two velocities.} The two-velocity step diffuses the mass-weighted
mixture of the transported intrinsic velocities with the operator
of~\cref{eq:viscosity:operator}, including its viscosity mixture, and gives
both phases the same increment~\thesis{\S7.5.2, step T3.II}:
\begin{equation}\label{eq:viscosity:two-velocity}
  \begin{gathered}
  \onevelo^{\stage{1}}_f=\frac{\dens^{\liq}\clampfrac{f}\,\velo^{\liq,\stage{1}}_f+\dens^{\gas}(1-\clampfrac{f})\,\velo^{\gas,\stage{1}}_f}{\stagdens_f},
  \qquad
  \onevelo^{\stage{2}}\ \text{from~\cref{eq:viscosity:phase-midpoint} with }\onevelo^{\stage{1}},\\
  \velo^{\phase,\stage{2}}_f=\velo^{\phase,\stage{1}}_f+\roundpar{\onevelo^{\stage{2}}_f-\onevelo^{\stage{1}}_f},
  \qquad \phase\in\{\liq,\gas\},
  \end{gathered}
\end{equation}
with the clamped staggered fraction~\cref{eq:forces:density}, a product with
an absent (\code{NaN}) phase taken as zero, $\onevelo^{\stage{1}}_f=0$ on
faces without open volume, and absent values left absent. The velocity jump
$\velo^{\gas}_f-\velo^{\liq}_f$ is unchanged and, on the faces with open
volume, the mixture of the new phase velocities is $\onevelo^{\stage{2}}$; on
a covered face the mixture is zero by definition, while
$\onevelo^{\stage{2}}$ is the value that the penalty rows of the covered
unknowns leave there (at most $2.6\times10^{-5}$ of the largest value in the
cut fixture; derivation audit, \code{proof_review/eq-viscosity-two-velocity/}).
The combined momentum
$\sum_f\opendualvol{f}\roundpar{\dens^{\liq}\clampfrac{f}\velo^{\liq}_f+\dens^{\gas}(1-\clampfrac{f})\velo^{\gas}_f}$,
with absent values counted as zero and $\opendualvol{f}=\dualvol{f}$ without
a solid, therefore changes exactly as the mixture momentum
$\sum_f\stagdens_f\opendualvol{f}\onevelo_f$, on a periodic mesh without
solid as in~\cref{eq:viscosity:momentum}, and a continuous velocity,
$\velo^{\liq}=\velo^{\gas}$, is diffused exactly as in the one-velocity
model. With \apitype{MimeticProjection} the phase values instead take the
implicit midpoint step of the same stress in the kinetic metric $\mathcal M$
of that projection~\cref{eq:projection:mimetic-viscous}
(\cref{sec:projection:mimetic}), which does not increase
$u^{\transpose}\mathcal Mu$ for any $\dt$; it equals
\cref{eq:viscosity:two-velocity} where $\mathcal M$ is the staggered phase
mass, and near the cut cells the jump and the combined momentum above are no
longer those of the common increment.

\emph{Interphase drag.} The common increment leaves the jump
$\jumpvelo_f=\velo^{\gas}_f-\velo^{\liq}_f$ of a face whose staggered volume
holds both phases undamped, although viscosity couples the phases across the
interface: a thin layer of one phase slips freely through the other (with a
per-phase surface-tension force the capillary wave at density ratio one has
the error $0.009$--$0.011$ instead of $0.0022$, audit note on small apertures,
\S5). With \apitype{InterphaseDrag}, the default of \apitype{TwoVelocityModel}
for both projections, every face with two finite phase values,
$0<\clampfrac{f}<1$ and an interface in its control volume adds a
Navier-slip interface shear $\pm\dragcond_f\jumpvelo_f/(1+\mathrm{Re}_f)$ to
the two phase momenta. With the staggered phase masses per volume
$\facemass{\liq}{f}=\dens^{\liq}\clampfrac{f}$,
$\facemass{\gas}{f}=\dens^{\gas}(1-\clampfrac{f})$ and
$\stagdens_f=\facemass{\liq}{f}+\facemass{\gas}{f}$,
\begin{equation}\label{eq:viscosity:drag}
  \begin{gathered}
  \dragcond_f=\frac{\visc_{\interface}(\clampfrac{f})\,\abs{\interface_f}\,w_f}{\delta_f},
  \qquad
  \lambda_f=\frac{\dragcond_f}{\stagvol{f}}\roundpar{\frac{1}{\facemass{\liq}{f}}+\frac{1}{\facemass{\gas}{f}}},
  \qquad
  \mathrm{Re}_f=a_f\abs{\jumpvelo_f},\quad a_f=\frac{h_f}{\nu_{\interface}},\quad
  \nu_{\interface}=\min\roundpar{\frac{\visc^{\liq}}{\dens^{\liq}},\frac{\visc^{\gas}}{\dens^{\gas}}},\\
  \frac{\dd\jumpvelo_f}{\dd\timevar}=-\frac{\lambda_f\jumpvelo_f}{1+a_f\abs{\jumpvelo_f}},
  \qquad
  a_f\abs{\jumpvelo^{+}_f}e^{a_f\abs{\jumpvelo^{+}_f}}=a_f\abs{\jumpvelo_f}e^{a_f\abs{\jumpvelo_f}}e^{-\lambda_f\dt},\\
  \velo^{\liq}_f\leftarrow\velo^{\liq}_f+\frac{\facemass{\gas}{f}}{\stagdens_f}\roundpar{\jumpvelo_f-\jumpvelo^{+}_f},
  \qquad
  \velo^{\gas}_f\leftarrow\velo^{\gas}_f-\frac{\facemass{\liq}{f}}{\stagdens_f}\roundpar{\jumpvelo_f-\jumpvelo^{+}_f},
  \end{gathered}
\end{equation}
the exact solution of
$\facemass{\liq}{f}\stagvol{f}\,\dot\velo^{\liq}_f=\dragcond_f\jumpvelo_f/(1+\mathrm{Re}_f)=-\facemass{\gas}{f}\stagvol{f}\,\dot\velo^{\gas}_f$
over the step (the Lambert-$W$ relation, solved by Newton's method from below,
where the iterates of the concave $w+\log w$ increase monotonically), applied
to the result of~\cref{eq:viscosity:two-velocity} with CMDGFM (with
\apitype{MimeticProjection} the same stress, with $\mathrm{Re}_f$ of the jump
at the start of the step, enters the implicit midpoint in the metric of that
projection,~\cref{eq:projection:mimetic-viscous-drag}, which keeps the
combined step non-expansive in that metric); $h_f$ is the mean cell
width along $\fnormal{f}$; $\nu_{\interface}$ depends on the phase properties
only (no volume fraction), and the layer counts as resolved only if it is
resolved in both phases (the effusivity-weighted diffusivity
$\roundpar{(\visc^{\liq}+\visc^{\gas})/(\sqrt{\dens^{\liq}\visc^{\liq}}+\sqrt{\dens^{\gas}\visc^{\gas}})}^2$
of two viscous half-spaces gave the same wave errors and a larger
Kelvin--Helmholtz error at density ratio ten). $\mathrm{Re}_f=\abs{\jumpvelo_f}h_f/\nu_{\interface}$ is the
cell Reynolds number of the jump: a grid-scale jump spreads by viscosity to the
thickness $\delta_s=h_f\mathrm{Re}_f^{-1/2}$ while it crosses a cell. Three
regimes follow from the slip length $\delta_f\mathrm{Re}_f$: a resolved shear
layer ($\delta_s\gg h$, $\mathrm{Re}_f\ll1$) is coupled without slip, the
exponential decay $e^{-\lambda_f\dt}$; an unresolved vortex sheet
($\delta_s\ll h$, $\mathrm{Re}_f\gg1$), which the two-velocity model
represents by the jump, slips freely; in between the stress is that of a
layer of thickness $\delta_f(1+\mathrm{Re}_f)$.
Here $\visc_{\interface}$ is the cell viscosity mixture of
\cref{sec:viscosity:stress} at the staggered fraction (for a flat interface
the series mean of the phases at the distances of their centroids from it)
and $\stagvol{f}$ the volume of the mixture solve. The control volume is the
one of the staggered fraction and the masses: each adjacent cell contributes
half the open part of the face's strip with the cell's phase fractions, so
its open-cell PLIC chord and phase centroids (the PLIC half-plane clipped to
the cell outside a solid's plane) enter scaled by that contribution over the
open cell area, and the liquid fraction of the volume is $\clampfrac{f}$ on
level transitions and cut cells as well. $\abs{\interface_f}$ is the sum of
the scaled chords, the interface normal $\inormal$ their length-weighted mean
normal, and $\delta_f$ the distance of the liquid and gas centroids of the
volume along $\inormal$, at least a quarter of the mean cell width along
$\fnormal{f}$. A flat interface parallel to the face gives
$\abs{\interface_f}=h$, $\delta_f=h/2$, so
$\lambda_f=2\nu/(h^2\clampfrac{f}(1-\clampfrac{f}))$ at equal kinematic
viscosities and densities; one normal to the face gives half the chord at
$\delta_f=h$, and between a pure liquid and a pure gas cell the face itself
is the interface. The default component weight (\code{component=:tangential})
is $w_f=\itangent_k^2/(\itangent_1^4+\itangent_2^4)$ for the face axis $k$:
a jump $J\itangent$ along the interface then gives
$\sum_k\itangent_k^2w_k=1$ per unit chord length for every orientation, as
$w_f=1$ (\code{component=:full}) does, while a jump $J\inormal$ normal to it
gives $2\itangent_1^2\itangent_2^2/(\itangent_1^4+\itangent_2^4)$, zero for a
grid-aligned interface and one at $45^\circ$, where the face components
cannot separate the two. Between a pure liquid and a pure gas cell the
tangential weight is zero: the face component is the normal velocity, and the
tangential velocities there sit on one-phase faces, coupled by the mixture
operator. The tangential and full weights reduce the capillary-wave error
with a per-phase surface-tension force at density ratio one from $0.0112$ to
$0.0019$ and $0.0023$ at $n=16$ and a quarter-period step factor $1/8$, and
otherwise differ by a few percent; the normal velocity jump is the
projection's. The slip law decides the rest: with the no-slip conductance
alone (no $\mathrm{Re}_f$) the viscous Kelvin--Helmholtz case
\code{r1_thin} with the mimetic projection loses its unresolved vortex sheet
and decays at $\lambda/h=64$ (growth-rate error $182\,\%$ instead of
$3.3\,\%$, $11.5\,\%$ instead of $3.3\,\%$ at $128$); the sheet's
$\mathrm{Re}_f$ is $6$--$1400$ (10--90\,\% quantiles of its mixed faces at
$\lambda/h=64$; $2.4$--$430$ at $128$), the slipping layer of the capillary
wave's $0.04$ (median; $0.67$ at 90\,\%), and the slip law gives $3.4\,\%$ and
$3.1\,\%$ with the same wave errors (these drag numbers were measured with
the exact update after the earlier common viscous increment).

With the mimetic projection the drag is a term of the implicit-midpoint step
in the metric~\cref{eq:projection:mimetic-viscous-drag}: the stress and the
drag together never increase $u^{\transpose}\mathcal Mu$, for every $\dt$.
(The exact update above, dissipative in the staggered masses, did positive
work in the metric in 97\,\% of the steps of \code{r10_thick}, $\nu^{\gas}/\nu^{\liq}=4$,
density ratio ten, and raised its growth-rate error from $16.5\,\%$/$21.5\,\%$
to $54\,\%$/$106\,\%$ at $\lambda/h=64/128$.) Since $\mathrm{Re}_f$ is taken
from the jump at the start of the step, the midpoint drag is second order in
time only where $\mathrm{Re}_f$ is fixed (a linear drag); with the lagged
$\mathrm{Re}_f$ it is first order. Measured with the metric step:
\code{r10_thick} gives $21.9\,\%$/$23.0\,\%$ against $0.05\,\%$/$22.9\,\%$
without drag and \code{r1_thick} $49.9\,\%$/$3.6\,\%$ against
$11.8\,\%$/$0.24\,\%$, while \code{r1_thin}, \code{r10_thin} and the
capillary-wave gain ($0.0015$--$0.0017$ instead of $0.008$--$0.014$ at
density ratio one) are kept. With CMDGFM the drag changes the results
by less than $10^{-4}$, because that projection recreates the tangential jump
of the cut faces in every step.

The exact update (CMDGFM) changes $\facemass{\liq}{f}\velo^{\liq}_f+\facemass{\gas}{f}\velo^{\gas}_f$
only by round-off, so the mixture momentum and the mixture
$\onevelo^{\stage{2}}$ are those of~\cref{eq:viscosity:two-velocity}; it
never increases $\abs{\jumpvelo_f}$, so it decreases the jump energy
$\half\stagvol{f}\redmass{f}\jumpvelo_f^2$ and leaves the mixture energy
unchanged; it vanishes for a zero jump (rest, a continuous velocity, the
static drop), on faces with one phase and for zero viscosities, and locks the
phases at fixed mixture momentum as $\lambda_f\dt\to\infty$. Since the common
increment leaves the jump unchanged and the drag leaves the mixture
unchanged, the two updates commute: a Strang split (half drag, mixture
solve, half drag) gives the same result as this sequence, and the scheme has
no splitting error between them. It is not the coupled solve of per-phase
stress operators with the interface rows $\dragcond_f$ added to
$\strainmat^{\transpose}\stressw\strainmat$, which the code does not
implement.

\begin{sloppypar}
\emph{Checks} (\code{midpoint_check}, \code{two_velocity_viscosity_check},
\code{single_phase_midpoint_check}). On the composite, cut and periodic
composite drop fixtures (densities $(1000,1)$, $\visc=(0.1,0.002)$,
$\dt=0.05$, tolerance $10^{-13}$), \cref{eq:viscosity:phase-midpoint} holds
to $\chkViscStrong$ (strong residual relative to the velocity),
\cref{eq:viscosity:energy} to $\chkViscEnergy$ of the dissipation, and on
the periodic mesh the momentum changes by $\chkViscMomentum$ (relative) while
$\strainmat^{\transpose}\stressw\strainmat\faceone{k}$ is $\chkViscConstant$ of
the largest stiffness entry. For two phase velocities that differ, the
mixture~\cref{eq:viscosity:two-velocity} equals an independent evaluation to
$\chkViscMixture$, the jump changes by $\chkViscJump$ and the mixture of the
updated phases equals $\onevelo^{\stage{2}}$ to $\chkViscRecombine$
(relative). The single-phase solve
satisfies~\cref{eq:viscosity:single-midpoint} to $\chkViscSingleMidpoint$ on
the walled composite mesh.
\end{sloppypar}

\begin{deviation}{Single-phase diffusion time level}\label{dev:viscosity:time-level}
The thesis' implicit midpoint step T3 acts on the transported velocity
$\velo^{\stage{1}}$ in both halves~\thesis{\S7.2, eq.~(7.52)}, which the
two-phase step~\cref{eq:viscosity:phase-midpoint} follows. The single-phase
step~\cref{eq:viscosity:single-midpoint} evaluates the explicit half at the
pre-transport velocity $\velo^{\timelevel{n}}$, which loses the energy
balance~\cref{eq:viscosity:energy} also where $\dualvolmat\lapstag{k}$ is
symmetric (it misses by $3\,\%$ in a uniform periodic test of the
derivation audit, \code{proof_review/eq-viscosity-single-midpoint/}).
\end{deviation}

\begin{implnote}
\apifun{implicit_midpoint_diffusion!}\code{(destination, convected, old_velocity, operator, dt, solver, work; time, tolerance, max_cycles)}~\cref{eq:viscosity:single-midpoint}
in \srcsym{src/navier_stokes/implicit_diffusion.jl}{implicit_midpoint_diffusion!}
with \srcsym{src/navier_stokes/implicit_diffusion.jl}{_add_wall_source!},
the shift \srcsym{src/navier_stokes/implicit_diffusion.jl}{_update_staggered_elliptic!}
of \apitype{StaggeredEllipticOperator}\code{(operator; reaction=1, diffusion=dt*nu/2)}
and \apitype{StaggeredEllipticSolver}; the two-phase
form~\cref{eq:viscosity:phase-midpoint} in
\srcsym{src/two_phase/diffusion/variable_diffusion_solver.jl}{implicit_midpoint_diffusion!}
with \srcsym{src/two_phase/diffusion/variable_diffusion_solver.jl}{VariableStaggeredDiffusionSolver};
stopping criterion \srcsym{src/solvers/multigrid/elliptic_solver.jl}{_solve_krylov_with_stats!};
stages \srcsym{src/single_phase/stepping.jl}{_diffusion_substep!},
\srcsym{src/two_phase/stepping.jl}{_diffusion_substep!} and
\srcsym{src/two_phase/two_velocity/stepping.jl}{_diffusion_substep!} with
\srcsym{src/two_phase/two_velocity/stepping.jl}{_mass_weighted_one_velocity!}
and \srcsym{src/two_phase/two_velocity/stepping.jl}{_apply_two_velocity_diffusion_update!}~\cref{eq:viscosity:two-velocity}
(\srcsym{src/two_phase/projection/mimetic_projection.jl}{_mimetic_viscous_step!} with the mimetic projection);
the drag~\cref{eq:viscosity:drag} \apitype{InterphaseDrag} in
\srcsym{src/two_phase/two_velocity/interphase_drag.jl}{_interphase_drag!} with
\srcsym{src/two_phase/two_velocity/interphase_drag.jl}{_interphase_drag_rate},
\srcsym{src/two_phase/two_velocity/interphase_drag.jl}{_interphase_contact},
\srcsym{src/two_phase/two_velocity/interphase_drag.jl}{_cell_phase_geometry} and
\srcsym{src/two_phase/two_velocity/interphase_drag.jl}{_clip_convex}; with the
mimetic projection
\srcsym{src/two_phase/two_velocity/interphase_drag.jl}{_interphase_conductance} and
\srcsym{src/two_phase/two_velocity/interphase_drag.jl}{_interphase_drag_coefficients!}.
Tests: \testsym{test/operators/staggered_diffusion.jl}{implicit midpoint staggered elliptic solve},
\testsym{test/operators/staggered_diffusion.jl}{dimensionless diffusion tolerance and failure report},
\testsym{test/navier_stokes/two_phase_navier_stokes.jl}{Two-phase diffusion uses the thesis split state}
(the step equals the direct solve of~\cref{eq:viscosity:phase-midpoint} about
the transported state),
\testsym{test/two_phase/variable_staggered_diffusion.jl}{Variable stress shifted solve on fixed and adapted trees},
\testsym{test/two_phase/two_velocity/diffusion.jl}{Two-velocity implicit diffusion identities},
\testsym{test/two_phase/two_velocity/diffusion.jl}{Two-velocity diffusion recovers the continuous-velocity limit},
\testsym{test/two_phase/two_velocity/interphase_drag.jl}{Interphase drag geometry},
\testsym{test/two_phase/two_velocity/interphase_drag.jl}{Interphase drag relaxes the jump at fixed mixture momentum},
\testsym{test/two_phase/two_velocity/interphase_drag.jl}{Interphase drag is inactive at rest, for equal phases and without viscosity},
\testsym{test/two_phase/two_velocity/interphase_drag.jl}{Interphase drag locks the phases as lambda dt grows},
\testsym{test/two_phase/two_velocity/interphase_drag.jl}{Interphase drag control volume on transitions, cut cells and small bubbles},
\testsym{test/two_phase/two_velocity/interphase_drag.jl}{Interphase drag semigroup and zero step},
\testsym{test/two_phase/two_velocity/interphase_drag.jl}{Interphase drag slip law},
\testsym{test/two_phase/two_velocity/interphase_drag.jl}{Two-velocity viscous step without drag is bitwise the mixture step},
\testsym{test/two_phase/two_velocity/interphase_drag.jl}{Interphase drag in the mimetic viscous step},
\testsym{test/navier_stokes/navier_stokes.jl}{single-phase material scaling}.
\end{implnote}
